\documentclass[12pt,a4paper]{article}

% Mathematics
\usepackage{amsmath,amssymb}
\usepackage{slashed}

% Graphics
\usepackage{graphicx}
\usepackage{epsfig}
\usepackage{epstopdf}

% Tables and page layout
\usepackage{multirow}
\usepackage{array}
\usepackage{needspace}
\newcolumntype{P}[1]{>{\raggedright\arraybackslash}p{#1}}
\usepackage{placeins}
\setlength{\parskip}{10pt}

% Text formatting and URLs
\usepackage{microtype}
\setlength{\emergencystretch}{3em}
\usepackage{url}

% Colors and cross-document citations
\usepackage[dvipsnames]{xcolor}

% Hyperlinks (load near the end)
\usepackage[
    colorlinks=true,
    citecolor=purple,
    linkcolor=blue
]{hyperref}

\begin{document}

\subsection*{Response to the Referee}

\noindent Dear Editor,

This paper considers an anomaly-free, renormalizable model of an additional Abelian
gauge boson \(A'\) with chiral couplings to muons and taus. In this model, vector-like
leptons and a scalar field \(S\) are introduced to construct the ultraviolet completion
of the axial \(U_A(1)_{L_{\mu-\tau}}\) gauge theory. The phenomenology of the \(A'\) in
the axial limit is studied, and the relevant experimental constraints are discussed.
The paper also discusses a possible future muon-collider probe as a complementary test of the chiral interaction structure.

The referee summarizes our work as follows:

\begin{quote}
\itshape
The authors construct an anomaly-free, renormalizable model in which an additional Abelian gauge field \(A'\) has chiral couplings to muons and taus. This model realizes a limit where the \(A'\) coupling to the \(\mu-\tau\) current is axial. The second goal of the paper is to examine the phenomenology of this model in the axial limit.

The paper presents results that are useful for the model building community. However, I have some significant reservations about key points that either should be included or otherwise justified in being excluded from the manuscript. These are explained below. Revisions for these major points (even if only to clarify why they are not issues) are required before the paper is at the level of a PRD manuscript.
\end{quote}

We sincerely thank the referee for the positive assessment of the model-building construction and for identifying issues that require further clarification or quantitative analysis. We respond to each point below.

\medskip
\noindent {\bf Major Points 1)}. \begin{quote}\itshape
Phenomenology of the \(S\) scalar. The second half of the paper reviews a set of experimental bounds on only an effective theory of the \(A'\) axially coupled to the \(L_{\mu-\tau}\) current. This leaves a major gap between these two parts of the paper: the authors show that their UV theory requires a scalar \(S\) that cannot be arbitrarily decoupled from the low-energy dynamics, e.g. from Eq. (34) which states that \(m_S^2 < 4\pi\left(m_{A'}/(2g_X)\right)^2\).\end{quote}

\noindent {\bf Reply.}
The original one-VEV \(S+\mathrm{VLL}\) construction contains a physical radial scalar \(s\). At tree level,
\[
m_{A'}=2g_Xv_s,\qquad
m_s^2=2\lambda_Sv_s^2,\qquad
y_{s\ell\ell}\simeq\frac{m_\ell}{v_s},
\quad \ell=\mu,\tau.
\]
Thus, the scalar mass and its leading charged-lepton couplings are tied to the UV symmetry-breaking scale and scalar self-coupling; \(s\) is not an arbitrary spectator that can simply be removed from the UV theory. The detailed scalar searches and complementary low-mass constraints are addressed under Major Point~1.1, the point-dependent scalar-mass range and perturbativity condition under Major Point~1.2, and the Higgs-portal mixing assumption under Major Point~1.3. Figure~2 presents the \(A'\) axial-current constraints for the stated scalar-suppressed benchmark; it does not imply that the radial scalar is absent from the UV completion.

\medskip
\Needspace{8\baselineskip}
\noindent {\bf Major Points 1.1)}. \begin{quote}\itshape
It is thus an odd choice to exclude the scalar from the phenomenology of the effective theory: if the \(S\) is unavoidable at phenomenologically relevant scales, shouldn't one of the main points of the paper be that the only known UV completion of an axial \(L_{\mu-\tau}\) gauge boson has a prediction for a scalar particle in a mass range that can be probed by colliders?\end{quote}

\noindent {\bf Reply.}
Yes. The radial scalar is a physical and potentially testable prediction of the UV completion. Its collider phenomenology must be assessed using scalar-specific production and decay hypotheses, with the UV branching fractions included when translating published limits.

\noindent\textbf{Direct collider searches and UV mapping.}

BaBar and Belle searched for
\[
e^+e^-\to\tau^+\tau^-s,\qquad
s\to e^+e^-,\,\mu^+\mu^-,
\]
over approximately \(0.04\,\mathrm{GeV}\lesssim m_s\lesssim6.5\text{--}7\,\mathrm{GeV}\)~\cite{BaBar:2020jma,Belle:2022gbl}. In the mass-proportional benchmark used in these analyses,
\[
\mathcal L\supset-\xi\sum_\ell\frac{m_\ell}{v_h}s\bar\ell\ell,
\qquad \xi=\frac{v_h}{v_s}.
\]
The Belle result \(\xi^{\rm lim}=0.483\) at \(m_s=1\,\mathrm{GeV}\)~\cite{Belle:2022gbl} therefore gives
\[
v_s>\frac{v_h}{\xi^{\rm lim}}\simeq509\,\mathrm{GeV},
\qquad
y_{s\mu\mu}<2.1\times10^{-4},
\]
and, using \(m_{A'}=2g_Xv_s\),
\[
g_X<9.8\times10^{-4}
\left(\frac{m_{A'}}{\mathrm{GeV}}\right)
\]
at \(90\%\) confidence level for this illustrative \(m_s=1\,\mathrm{GeV}\) point with \(m_{A'}>m_s/2\), so that the two-body channel \(s\to A'A'\) is closed. This example illustrates the mass-dependent scalar-search translation; it is separate from the \(m_s\gtrsim10\,\mathrm{GeV}\) scalar-suppressed reference benchmark used for Fig.~2. If \(m_{A'}<m_s/2\), the \(A'A'\) mode is open and the visible-channel limit must be interpreted with the UV branching fraction rather than applied as the same \(g_X\) bound. For the \(\mu\)- and \(\tau\)-lepton couplings relevant to the associated production and dilepton/ditau channels, the mass-proportional benchmark has the same Yukawa normalization as the radial scalar in our construction:
\[
y_{s\ell\ell}=\frac{m_\ell}{v_s}
=\frac{2g_Xm_\ell}{m_{A'}}
=\xi\frac{m_\ell}{v_h},
\qquad
\xi=\frac{v_h}{v_s},
\qquad \ell=\mu,\tau.
\]
At fixed \((m_{A'},g_X)\), this Yukawa coupling is independent of \(m_s\); changing \(m_s\) changes the scalar self-coupling and the kinematically open decay channels, not the mass-proportional Yukawa normalization. The quoted Belle example at \(m_s=1\,\mathrm{GeV}\) uses the dimuon channel. In the minimal construction with negligible portal mixing, the tree-level radial-scalar Yukawa couplings are those to \(\mu\) and \(\tau\); the generic \(s\to e^+e^-\) channel in the experimental benchmark is therefore not used as a direct constraint on this UV benchmark.

For \(m_s>2m_{A'}\), the gauge completion predicts the two-vector partial width
\[
\Gamma(s\to A'A')
=\frac{m_s^3}{32\pi v_s^2}
\left(1-4x+12x^2\right)\sqrt{1-4x},
\qquad
x=\frac{m_{A'}^2}{m_s^2},
\]
while each kinematically open charged-lepton channel has
\[
\Gamma(s\to\ell^+\ell^-)
=\frac{m_\ell^2m_s}{8\pi v_s^2}
\left(1-\frac{4m_\ell^2}{m_s^2}\right)^{3/2},
\qquad \ell=\mu,\tau.
\]
When both channels are open, their ratio is
\[
\frac{\Gamma(s\to A'A')}{\Gamma(s\to\ell^+\ell^-)}
=
\frac{m_s^2}{4m_\ell^2}
\frac{(1-4x+12x^2)\sqrt{1-4x}}
{\left(1-4m_\ell^2/m_s^2\right)^{3/2}}.
\]
Away from the \(A'A'\) threshold this ratio can be large, so the \(A'A'\) mode can dominate over an individual charged-lepton partial width. The visible branching fraction is consequently
\[
\begin{aligned}
\operatorname{BR}(s\to\ell^+\ell^-)
&=\frac{\Gamma(s\to\ell^+\ell^-)}{\Gamma_s^{\mathrm{tot}}},\\
\Gamma_s^{\mathrm{tot}}&\supset
\Theta(m_s-2m_\mu)\Gamma(s\to\mu^+\mu^-)\\
&\quad+\Theta(m_s-2m_\tau)\Gamma(s\to\tau^+\tau^-)\\
&\quad+\Theta(m_s-2m_{A'})\Gamma(s\to A'A').
\end{aligned}
\]
Thus the event yield in a dilepton search depends on the production rate, this model-dependent visible branching fraction, and the corresponding selection efficiency. When \(s\to A'A'\) is open, the resulting \(A'\) decays can also change the final-state topology. The published B-factory limits therefore constrain their stated mass-proportional visible-decay benchmark. Their direct translation to the present UV parameter plane applies where the UV branching fractions reproduce the assumed visible benchmark; when \(s\to A'A'\) is open, the interpretation instead requires the model-dependent visible branching fractions and the corresponding final-state selection.

\noindent\textbf{Belle~II scalar-specific searches.}

Belle~II provides scalar-specific searches; its published resonance results in this mass range are not uniformly spin-one-only. The \(e^+e^-\to\mu^+\mu^-X\), \(X\to\tau^+\tau^-\) analysis covers \(3.6\text{--}10\,\mathrm{GeV}\) and tests a Higgs-like spin-zero leptophilic scalar with mass-proportional charged-lepton couplings, matching the \(\mu\)- and \(\tau\)-Yukawa normalization of our radial scalar~\cite{Belle-II:2023ydz}. The four-muon analysis covers \(0.212\text{--}9\,\mathrm{GeV}\) and includes a muonphilic-scalar hypothesis alongside the \(L_\mu-L_\tau\) vector boson; this is complementary but is not identical to the full mass-proportional scalar pattern of our UV construction~\cite{Belle-II:2024wtd}. When \(s\to A'A'\) is open, it changes the total width and visible \(\tau^+\tau^-\) branching fraction, and any cascade topology requires its corresponding efficiency. The published Belle~II scalar limits therefore provide relevant sensitivity, while their translation to the radial scalar requires the UV branching fractions and event selection. General singlet--Higgs mixing and additional-scalar searches provide complementary context~\cite{Robens:2015singlet,Robens:2020twoSinglets,Boiarska:2019scalarPortal}.

\noindent\textbf{Complementary low-mass constraints and fifth-force matching.}

Complementary low-mass constraints on the scalar sector are summarized in Table~II. Direct scalar-mediator limits are expressed in terms of \(y_{s\mu\mu}\), which is related to the UV parameters through
\[
y_{s\mu\mu}=\frac{2g_Xm_\mu}{m_{A'}}.
\]
Fifth-force and photon-mediated gamma-ray limits instead rely on loop-induced couplings and retain the source-specific assumptions stated below.

The binary-pulsar entry uses Liu and Tang's scalar-mediator result, whose coupling has the same Yukawa normalization as \(y_{s\mu\mu}\). We quote the \(\mathcal O(10^{-21})\) best-sensitivity point separately from the reported mass interval; it is not a uniform limit across that interval. The \(L_\mu-L_\tau\) vector-mediator result of Dror et al. is a distinct constraint~\cite{Liu:2025zuz,Dror:2019uea}.

\begin{center}
\small
\begin{tabular}{P{0.26\linewidth}P{0.36\linewidth}P{0.25\linewidth}}
\hline
Scalar-mass range & Published constraint & Source-reported bound / applicability\\
\hline
\(10^{-17}\,\mathrm{eV}\lesssim m_s\lesssim2\times10^{-15}\,\mathrm{eV}\) & Binary-pulsar periastron advance (95\% CL), ultralight scalar~\cite{Liu:2025zuz} & \(\mathcal O(10^{-21})\) at the best-sensitivity point \(m_s\simeq4\times10^{-16}\,\mathrm{eV}\); not a uniform bound over the interval\\
\(m_s\lesssim30\,\mathrm{meV}\) (crossover scale, not a hard endpoint) & Fifth-force tests from the muon-only loop-induced scalar--proton coupling~\cite{Fiorillo:2026wso,Hoskins:1985tn,Kapner:2006si,Blakemore:2021zna} & Mass-dependent source curve; \(\sim30\) meV is the crossover below which it overtakes astrophysical limits; no universal direct \(y_{s\mu\mu}\) limit\\
\(m_s\lesssim100\,\mathrm{keV}\) & Neutron-star cooling~\cite{Fiorillo:2026wso} & \(\mathcal O(10^{-13})\); Table~1 gives \(7.4\times10^{-14}\) for \(m_s\lesssim1\,\mathrm{keV}\)\\
\(m_s\lesssim\mathcal O(10)\,\mathrm{MeV}\) & SN~1987A neutrino energy loss, conservative cold-SFHo free-streaming benchmark~\cite{Caputo:2021rux,Fiorillo:2026wso} & \(\lesssim4.2\times10^{-9}\) over the mass range shown in Caputo et al.'s Fig.~2\\
\(m_s\lesssim\mathcal O(10)\,\mathrm{MeV}\) & SN~1987A \(\gamma\)-ray fluence, cold-SN benchmark~\cite{Caputo:2021rux} & \(g_{\phi\mu}^{\rm source}<1.11\times10^{-10}\sqrt{\mathrm{MeV}/m_s}\); mass dependent and stronger than neutrino cooling for \(m_s\gtrsim0.70\,\mathrm{keV}\)\\
\(0.1\,\mathrm{MeV}\lesssim m_s\lesssim\mathcal O(10)\,\mathrm{MeV}\) & Diffuse supernova \(\gamma\)-ray background, short-lived scalar branch~\cite{Caputo:2021rux} & \(g_{\phi\mu}^{\rm source}\lesssim4.2\times10^{-11}\) for \(n_7^{\rm cc}=0.58\); the long-lived branch is mass dependent\\
\hline
\end{tabular}
\end{center}

The entries in Table~II refer to distinct observables and source assumptions and are not combined into a single exclusion band. Two mass-range qualifications are important. Fiorillo et al.'s v2 neutron-star cooling result is mass dependent up to approximately \(100\,\mathrm{keV}\); the more specific \(7.4\times10^{-14}\) value in their Table~1 applies only for \(m_s\lesssim1\,\mathrm{keV}\) and must not be read as a constant bound over the full mass interval~\cite{Fiorillo:2026wso}. For SN~1987A neutrino energy loss, Fiorillo et al.'s Table~1 quotes \(4.2\times10^{-9}\) for its \(m_s\lesssim1\,\mathrm{keV}\) benchmark, whereas Caputo et al.'s Fig.~2 displays the conservative cold-SFHo free-streaming boundary out to \(m_s\sim\mathcal O(10)\,\mathrm{MeV}\), remaining approximately at that level over the plotted range. We use the mass coverage of the latter curve, not the mass label of the former tabulated benchmark. This is a source-specific free-streaming cooling criterion, not a complete trapping or explosion-energy exclusion~\cite{Caputo:2021rux,Fiorillo:2026wso}. The SN~1987A gamma-ray fluence and diffuse-supernova gamma-ray background are separate observables, listed separately in Table~II. The quoted photon-mediated results use Caputo et al.'s source normalization \(g_{\phi\mu}^{\rm source}\) and muon-only treatment; their production and loop approximations are restricted to \(m_\phi\lesssim10\,\mathrm{MeV}\)~\cite{Caputo:2021rux}.

The fifth-force result is an indirect constraint. In Fiorillo et al.'s low-momentum, muon-only benchmark, the muon triangle induces a scalar--photon interaction and an effective proton coupling,
\[
g_{\phi p}^{(\mu)}=\frac{\alpha^2}{4}g_{\phi\mu}.
\]
Ordinary-matter fifth-force experiments constrain this loop-induced proton coupling, rather than a tree-level scalar--muon Yukawa vertex. The mass-dependent curve in Fig.~1 of Fiorillo et al. overtakes the astrophysical limits below an approximate crossover scale of \(m_\phi\sim30\,\mathrm{meV}\); this is not a sharp upper endpoint of experimental sensitivity~\cite{Fiorillo:2026wso}.

In the one-VEV UV construction, the radial scalar couples to both charged leptons,
\[
y_{s\mu\mu}=\frac{m_\mu}{v_s},\qquad
y_{s\tau\tau}=\frac{m_\tau}{v_s},\qquad
\frac{y_{s\mu\mu}}{m_\mu}
=\frac{y_{s\tau\tau}}{m_\tau}
=\frac{1}{v_s}.
\]
Consequently, in the fifth-force regime \(m_s\lesssim30\,\mathrm{meV}\ll2m_\mu\), the low-momentum electromagnetic form factor receives equal-sign contributions from both loops:
\[
\mathcal F_{s\gamma\gamma}(0)
\propto
\sum_{\ell=\mu,\tau}Q_\ell^2
\frac{y_{s\ell\ell}}{m_\ell}
=\frac{2}{v_s}.
\]
The contributions add because the amplitude is weighted by \(Q_\ell^2\), so the opposite \(U(1)_{L_\mu-L_\tau}\) charges do not cause a cancellation. Fiorillo et al.'s curve is therefore a muon-only source benchmark, not the complete fifth-force prediction of this UV model; translating it requires matching both the muon and tau loops.

For application of the source-level limits to the UV scalar, the decay \(s\to A'A'\) must be included whenever \(m_s>2m_{A'}\). When open, it contributes to the total width and can change the lifetime, visible branching fractions, cascade topology, and energy deposition relevant to a given search or astrophysical observable. The entries in Table~II are therefore source-level constraints; their UV interpretation depends on the corresponding production and decay assumptions.

For a source result that directly constrains the same scalar--muon Yukawa normalization, the one-VEV relation given above yields the conditional translation
\[
g_X\lesssim\frac{y_{s\mu\mu}^{\rm lim}(m_s)m_{A'}}{2m_\mu},
\]
at a specified \(m_s\) within the source's mass range. This is not a model-independent bound on \(g_X\). In particular, the fifth-force curve cannot be inserted into this expression as a direct Yukawa limit; its UV translation first requires the two-lepton loop matching.

\noindent\textbf{Scalar contribution to the muon anomalous magnetic moment.}

The scalar contribution to the muon anomalous magnetic moment is included together with the axial-vector contribution,
\[
\Delta a_\mu=\Delta a_\mu^{A'}+\Delta a_\mu^s,
\qquad
\Delta a_\mu^s=
\frac{y_{s\mu\mu}^2}{8\pi^2}
\int_0^1 dx\,
\frac{x^2(2-x)m_\mu^2}
{m_s^2(1-x)+x^2m_\mu^2}.
\]
For the reference regime \(m_s\gtrsim10\,\mathrm{GeV}\), the scalar term is below \(0.2\%\) of the pure-axial \(A'\) contribution. For lighter \(s\), the positive scalar term can instead compete with, or partially cancel, the negative axial contribution. Accordingly, the \((g-2)_\mu\) contour in Fig.~2 is evaluated for the stated scalar-suppressed benchmark, rather than being applied to arbitrary scalar masses.

Taken together, the direct searches and the complementary low-mass constraints establish that the radial scalar has its own phenomenology. Their interpretation is conditional on the production and decay structure of the one-VEV UV model; in particular, a published visible-channel limit is not automatically a model-independent exclusion when \(s\to A'A'\) changes the total width and branching fractions. The \(A'\) axial-current contours in Fig.~2 are therefore a separate, explicitly stated benchmark rather than a claim that scalar effects are absent.

\medskip

\noindent {\bf Major Points 1.2)}. \begin{quote}\itshape
What is the range of phenomenologically allowed \(S\) mass including perturbativity limits on \(\lambda\) and a motivated range of \(m_{A'}\) and \(g_X\). For example: perhaps the authors imagine a range of \(m_{A'}\) below some benchmark mass where certain experimental constraints on \(g_X\) become very strong. Then with that set of constraints (which should be stated clearly, see minor point \#1 below), one may give a range of \(m_S\) and comment on the experimental constraints.\end{quote}

\noindent {\bf Reply.}
The manuscript specifies the scalar-mass range point by point through the UV relations, rather than assigning a single scalar mass interval independently of \((m_{A'},g_X)\). In the one-VEV construction,
\[
v_s=\frac{m_{A'}}{2g_X},
\qquad
\lambda_S=\frac{m_s^2}{2v_s^2}
=\frac{2g_X^2m_s^2}{m_{A'}^2}.
\]
For the scalar-suppressed reference benchmark \(m_s\gtrsim10\,\mathrm{GeV}\) and the perturbativity criterion \(\lambda_S\leq4\pi\), the corresponding point-dependent range is
\[
10\,\mathrm{GeV}\lesssim m_s\leq\sqrt{8\pi}\,v_s
=\sqrt{2\pi}\,\frac{m_{A'}}{g_X}.
\]
It requires
\[
v_s\gtrsim\frac{10\,\mathrm{GeV}}{\sqrt{8\pi}}
\simeq1.995\,\mathrm{GeV},
\qquad
m_{A'}\gtrsim(3.99\,\mathrm{GeV})\,g_X.
\]
This is the theoretical consistency boundary represented by the \(v_s\simeq2\,\mathrm{GeV}\) line in Fig.~2. At each plotted point, \(v_s\) is derived and the scalar self-coupling is checked against the perturbativity criterion.

The lower-mass scalar regimes and the direct BaBar/Belle/Belle~II searches are discussed in detail under Major Point~1.1. Thus, the manuscript distinguishes the point-dependent perturbative mass range used for the \(A'\) benchmark from the mass-dependent scalar phenomenology. The dashed line is a theoretical consistency condition for the specified \(m_s\gtrsim10\,\mathrm{GeV}\) benchmark, not an experimental exclusion or a universal boundary for every possible scalar-sector choice.

\medskip
\noindent {\bf Major Points 1.3)}. \begin{quote}\itshape
Above (34): the authors ``assume negligible Higgs-portal mixing.'' Is this interaction not generated at loop level? It seems that \(S\)-\(h\) mixing cannot be avoided due to \(\mu\) and \(\tau\) loops. In turn, this seems like it would lead to strong experimental bounds on the UV completion. In fact, once there is mixing, is there not a resonance where the mixing becomes large as \(m_S\to m_h\)? It would be useful to quantify the size of this mixing over the allowed range of \(m_S\). Existing Higgs measurements may set the strongest bounds on the UV completion. If these bounds are obviously negligible (though again, I expect this to be strongly constrained as \(m_S\to m_h\)), then the authors should present this as the justification for ignoring \(m_S\) in the low-energy theory. (Though I am still skeptical that the \(S\) can be ignored since it cannot formally decouple.)\end{quote}

\noindent {\bf Reply.}
We agree that a vanishing Higgs-portal coupling is not enforced by the gauge symmetry: the operator \((H^\dagger H)(S^\dagger S)\) is allowed and contributes \((M_{\mathrm{loc}}^2)_{hs}=\lambda_{HS}v_hv_s\) after symmetry breaking, so the portal boundary condition must be stated at a matching scale and in a specified renormalization convention. The negligible-portal assumption is therefore used only as a conditional phenomenological benchmark, not as a symmetry-protected statement or a claim that the portal vanishes at every scale.

Generic singlet--Higgs mixing is constrained by the \(125\) GeV Higgs signal strengths and direct searches for additional scalar states; representative analyses are cited above. For the original one-singlet \(S+\mathrm{VLL}\) UV model considered here, we complement these general results with the model-specific VLL threshold matching and the finite-momentum \(h\)-\(s\) self-energy, evaluated using the full charged-lepton/VLL mass-eigenstate matrices as detailed below.

In the \(\overline{\mathrm{MS}}\) scheme, with \(\mathcal L\supset -V\) and \(V\supset\lambda_{HS}(H^\dagger H)(S^\dagger S)\), the zero-momentum effective-potential matching for the vector-like-lepton (VLL) sector gives
\[
\lambda_{HS}^{\mathrm{EFT}}(\mu_R)
=\lambda_{HS}^{\mathrm{UV}}(\mu_R)
-\sum_{\ell=\mu,\tau}
\frac{|y_\ell|^2|\lambda_\ell|^2}{8\pi^2}
+\cdots,
\]
where \(y_\ell\) and \(\lambda_\ell\) are the Yukawa couplings in the VLL interaction. This is a conditional zero-momentum threshold relation; it does not determine the finite-momentum off-diagonal self-energy or a physical scalar mixing angle.

To make the portal assumption explicit, we have evaluated the leading VLL threshold at the matching scale \(\mu_M=M_E\),
\[
\Delta\lambda_{HS}^{E}
=-\frac{1}{8\pi^2}\sum_{\ell=\mu,\tau}|y_\ell|^2|\lambda_\ell|^2,
\qquad
\lambda_{HS}^{\mathrm{EFT}}(\mu_M)
=\lambda_{HS}^{\mathrm{UV}}(\mu_M)+\Delta\lambda_{HS}^{E}.
\]
We use \(\lambda_{HS}^{\mathrm{EFT}}(M_E)=0\) as an explicit small-portal reference condition, which requires a nonzero UV portal \(\lambda_{HS}^{\mathrm{UV}}=-\Delta\lambda_{HS}^{E}\). This condition is a benchmark choice, not a symmetry prediction. In the leading charged-lepton/VLL mass-mixing approximation,
\[
\theta_{\ell L}\simeq\frac{y_\ell v_h}{\sqrt2M_E},
\qquad
\theta_{\ell R}\simeq\frac{\lambda_\ell v_s}{\sqrt2M_E},
\qquad
m_\ell\simeq\theta_{\ell L}\theta_{\ell R}M_E
\simeq\frac{y_\ell\lambda_\ell v_hv_s}{2M_E}.
\]
For illustration, taking \(M_E=20\,\mathrm{TeV}\), \(\theta_{\mu L}=\theta_{\tau L}=10^{-2}\), and the leading charged-lepton/VLL mass-mixing relations gives \(\Delta\lambda_{HS}^{E}\simeq-0.2108\) at \((m_{A'},g_X,v_s)=(10\,\mathrm{GeV},0.0705,70.9\,\mathrm{GeV})\), and \(\Delta\lambda_{HS}^{E}\simeq-0.01067\) at \((63\,\mathrm{GeV},0.1,315\,\mathrm{GeV})\). For these illustrative points only, we additionally impose \(|\lambda_\ell|\leq\sqrt{4\pi}\); this is a theory-consistency benchmark rather than a new exclusion contour in Fig.~2. The comparison demonstrates that setting the UV portal to zero is not equivalent to setting the matched EFT portal to zero, and that the latter must be stated as an input if the small-portal benchmark is retained.

The condition \(\lambda_{HS}^{\mathrm{EFT}}(M_E)=0\) removes the local dimension-four mixed mass insertion from the matched EFT at the matching scale; it does not set the UV portal to zero. In particular, \(\lambda_{HS}^{\mathrm{UV}}(M_E)=-\Delta\lambda_{HS}^{E}\) is nonzero in both benchmarks. The CP-even radial-scalar bubbles generated by this portal can be organized using
\[
g_{hhh}=6\lambda_Hv_h,\qquad
g_{sss}=6\lambda_Sv_s,\qquad
g_{hhs}=\lambda_{HS}^{\mathrm{UV}}v_s,\qquad
g_{hss}=\lambda_{HS}^{\mathrm{UV}}v_h.
\]
Their zero-momentum-subtracted contribution has the form
\[
\begin{aligned}
\Delta\Pi_{hs}^{(h,s)\,\mathrm{subset}}
&=\frac{\lambda_{HS}^{\mathrm{UV}}v_hv_s}{32\pi^2}
\Big[
6\lambda_H\Delta\bar B_{hh}
+6\lambda_S\Delta\bar B_{ss}\\
&\qquad
+2\lambda_{HS}^{\mathrm{UV}}\Delta\bar B_{hs}
\Big],\\
\Delta\bar B_{ij}
&\equiv\Delta\bar B_0(m_h^2;m_i^2,m_j^2).
\end{aligned}
\]
where the finite, subtraction-scale-independent function is
\[
\Delta\bar B_0(s;a,b)
=-\int_0^1 dx\,
\ln\!\left[
1-\frac{x(1-x)s}{xa+(1-x)b}
\right],
\qquad a=m_i^2,\quad b=m_j^2.
\]
For equal internal masses, this integral has the elementary closed form
\[
\boxed{
\Delta\bar B_0(s;a,a)
=2-2\sqrt{\frac{4a}{s}-1}\,
\arctan\!\left(\frac{1}{\sqrt{4a/s-1}}\right),
\qquad 0<s<4a.
}
\]
For unequal internal masses, define
\[
\begin{aligned}
\Lambda(s,a,b)&=(s-a-b)^2-4ab,\\
q&=\frac{\sqrt{-\Lambda(s,a,b)}}{2s},\\
t_0&=\frac{a-b-s}{2s},\qquad t_1=t_0+1.
\end{aligned}
\]
and
\[
\begin{aligned}
G(t)&=t\ln\!\left[\frac{s(t^2+q^2)}{\mu_0^2}\right]
-2t+2q\arctan\!\left(\frac{t}{q}\right),\\
I_0(a,b;\mu_0)&=
\frac{a\ln(a/\mu_0^2)-b\ln(b/\mu_0^2)}{a-b}-1.
\end{aligned}
\]
The mixed-mass result is then
\[
\boxed{
\Delta\bar B_0(s;a,b)
=I_0(a,b;\mu_0)-\big[G(t_1)-G(t_0)\big],
\qquad a\ne b,\quad \Lambda(s,a,b)<0.
}
\]
The arbitrary reference scale \(\mu_0\) cancels from this difference. All benchmark evaluations below lie on this real sub-threshold branch. Thus the \(hh\) and \(ss\) terms use the equal-mass expression, while the \(hs\) term uses the mixed-mass expression. At \(m_s=100\,\mathrm{GeV}\), the function triplet is \((\Delta\bar B_{hh},\Delta\bar B_{ss},\Delta\bar B_{hs})=(0.1862006,0.3150601,0.2366921)\); at \(m_s=m_h\), all three equal \(0.1862006\). The closed forms agree with independent 8,000-interval Simpson integration at relative differences below \(10^{-14}\).
Because \(\lambda_{HS}^{\mathrm{UV}}\) is fixed at the one-loop VLL-threshold order, the terms with one portal insertion are next order in this combined counting; the term with two portal insertions is one order further. In an \(R_\xi\) gauge, the corresponding Goldstone vertices are \(g_{hG_HG_H}=2\lambda_Hv_h\), \(g_{sG_HG_H}=\lambda_{HS}^{\mathrm{UV}}v_s\), \(g_{haa}=\lambda_{HS}^{\mathrm{UV}}v_h\), and \(g_{saa}=2\lambda_Sv_s\). A physical pole calculation must combine these with the gauge/ghost contributions and a consistent tadpole/VEV prescription. Tadpole and VEV-counterterm terms are local and cancel in the displayed momentum difference, but are required in a complete un-subtracted mass/pole renormalization. Thus this expression fixes the perturbative order and diagram structure; it is not by itself a complete two-loop UV--EFT matched self-energy. We do not add this isolated subset to \(\Pi_{hs}^{\mathrm{ret}}\) or use it to classify the near-degenerate region. In the gauge basis the singlet radial mode has no direct \(sWW\) or \(s\bar t t\) vertex, so direct one-loop \(W\)- or top-quark contributions to \(\Pi_{hs}\) are absent at the retained order; portal-induced effects in those loops enter at higher order.

For each charged-lepton flavour, the field-dependent mass matrix and gauge-basis scalar vertices are
\[
\begin{aligned}
\mathcal M_\ell(h,s)&=
\begin{pmatrix}
0&y_\ell(v_h+h)/\sqrt2\\
\lambda_\ell(v_s+s)/\sqrt2&M_{E_\ell}
\end{pmatrix},\\[3pt]
\Gamma_h^\ell&=
\begin{pmatrix}0&y_\ell/\sqrt2\\0&0\end{pmatrix},
\qquad
\Gamma_s^\ell=
\begin{pmatrix}0&0\\\lambda_\ell/\sqrt2&0\end{pmatrix}.
\end{aligned}
\]
For the \(\tau\) interaction the original model contains \(S^*\); its radial real-\(s\) vertex is the same for the real benchmark couplings. With
\[
U_L^T\mathcal M_\ell(0,0)U_R
=\operatorname{diag}(m_{\ell1},m_{\ell2}),
\qquad
G_a^\ell=U_L^T\Gamma_a^\ell U_R\quad(a=h,s),
\]
all four internal mass-eigenstate pairs are included for both \(\mu\) and \(\tau\). In the following kernels the flavour label \(\ell\) on \(G_a^\ell\) is suppressed.

Define
\[
C_{ab}^{ij}=(G_a)_{ij}(G_b)_{ij}^{*}
 +(G_a)_{ji}^{*}(G_b)_{ji},
\qquad
D_{ab}^{ij}=(G_a)_{ij}(G_b)_{ji}
 +(G_a)_{ji}^{*}(G_b)_{ij}^{*},
\]
and
\[
K_{ij}=\bar A_0(m_i^2)+\bar A_0(m_j^2)
 +(m_i^2+m_j^2-p^2)\bar B_0(p^2;m_i^2,m_j^2).
\]
The one-loop fermion self-energy in the source convention is
\[
\Pi_{ab}^{F,\ell}(p^2)
=-\frac{1}{16\pi^2}\sum_{i,j=1}^{2}
\left[C_{ab}^{ij}K_{ij}
+2m_im_jD_{ab}^{ij}\bar B_0(p^2;m_i^2,m_j^2)\right].
\]
We use \(\Gamma_{\rm inv}=p^2\mathbf1-\mathcal M_{\rm loc}^2+\Pi_{\rm RC}\) and \(\Pi_B=-\Pi_{\rm RC}\). The finite functions at \(\mu_R=M_E\) are
\[
\begin{aligned}
\bar A_0(m^2)&=m^2\left(1-\ln\frac{m^2}{\mu_R^2}\right),\\
\bar B_0(p^2;m_i^2,m_j^2)
&=-\int_0^1dx\,
\ln\!\left[
\frac{xm_i^2+(1-x)m_j^2-x(1-x)p^2-i0}{\mu_R^2}
\right].
\end{aligned}
\]

The full-theory self-energy contains the local portal-threshold constant and the light-lepton contribution reproduced in the EFT. The momentum-dependent hard term is isolated by
\[
\Delta\Pi_{\mathrm{hard}}(m_h^2)=
[\Pi_B^{\mathrm{full}}(m_h^2)-\Pi_B^{\mathrm{full}}(0)]
-[\Pi_B^{\mathrm{light,exact}}(m_h^2)-\Pi_B^{\mathrm{light,exact}}(0)],
\]
where the light-light term uses the exact light-state Yukawa couplings from the same mass diagonalization. The retained EFT soft contribution is the light-lepton bubble plus the dimension-five contact tadpole,
\[
\begin{aligned}
\Pi_{\mathrm{soft,bub+ct}}^B(p^2)
&=\sum_{\ell=\mu,\tau}\frac{q_\ell}{16\pi^2}
\left[4\bar A_0(m_\ell^2)
+(4m_\ell^2-p^2)\bar B_0(p^2;m_\ell^2,m_\ell^2)\right],\\
q_\ell&=\frac{2m_\ell^2}{v_hv_s}.
\end{aligned}
\]
We denote their sum by
\[
\Pi_{hs}^{\mathrm{ret}}(m_h^2)
=\Pi_{\mathrm{soft,bub+ct}}^B(m_h^2)
+\Delta\Pi_{\mathrm{hard}}(m_h^2).
\]
The soft term here is the specified bubble-plus-contact subset, not the complete set of all same-order EFT soft insertions and counterterms. Thus \(\Pi_{hs}^{\mathrm{ret}}\) is a conditional Yukawa-sector benchmark estimate rather than a fully closed model-level pole prediction.

For \(M_E=20\,\mathrm{TeV}\), \(\theta_{\mu L}=\theta_{\tau L}=10^{-2}\), \(m_h=125.25\,\mathrm{GeV}\), \(v_h=246.22\,\mathrm{GeV}\), and \(\mu_R=M_E\), the numerical decomposition is
\[
\resizebox{\textwidth}{!}{$
\begin{array}{c|cc}
[\mathrm{GeV}^2]&\text{Benchmark A}&\text{Benchmark B}\\ \hline
\Pi_{\mathrm{soft,bub+ct}}^B(m_h^2)&
-0.437308303-0.113173709i&
-0.098410275-0.025468201i\\
\Delta\Pi_{\mathrm{hard}}(m_h^2)&
+0.024030444&
+0.005407724\\
\Pi_{hs}^{\mathrm{ret}}(m_h^2)&
-0.413277859-0.113173709i&
-0.093002551-0.025468201i
\end{array}
$}
\]
The exact hard momentum difference agrees with the leading source-aligned derivative result to \(6.11\times10^{-4}\) relative at both benchmarks, and the zero-momentum hard comparator agrees with the leading portal-threshold estimate to \(1.79\times10^{-4}\) relative. These checks validate the stated fermionic subcalculation and its numerical implementation, but do not close the remaining same-order EFT soft/counterterm sectors.

For a clearly non-degenerate point in Benchmark B, we take \(m_s=100\,\mathrm{GeV}\). Here \(m_s\) denotes the unmixed diagonal singlet-mass input in the inverse propagator, not an already diagonalized physical pole mass, and
\[
\lambda_S=\frac{m_s^2}{2v_s^2}=0.05039.
\]
The singlet width in the mass interval considered is
\[
\begin{aligned}
\Gamma_s={}&
\Theta(m_s-2m_{A'})
\frac{m_s^3}{32\pi v_s^2}
(1-4x+12x^2)\sqrt{1-4x}\\
&+\sum_{\ell=\mu,\tau}
\Theta(m_s-2m_\ell)
\frac{m_\ell^2m_s}{8\pi v_s^2}
\left(1-\frac{4m_\ell^2}{m_s^2}\right)^{3/2},
\qquad x=\frac{m_{A'}^2}{m_s^2}.
\end{aligned}
\]
At this point \(s\to A'A'\) is closed because \(m_s<2m_{A'}=126\,\mathrm{GeV}\), \(h\to A'A'\) is closed because \(m_h<2m_{A'}\), and \(h\to ss\) is closed because \(m_s>m_h/2\). The width is \(\Gamma_s=0.126811\,\mathrm{MeV}\). The first-pass complex-pole solution gives a singlet-like pole at \(99.999999997\,\mathrm{GeV}\) with width \(0.126815\,\mathrm{MeV}\), and a Higgs-like pole at \(125.250000022\,\mathrm{GeV}\) with width \(4.069993\,\mathrm{MeV}\).

For this non-degenerate point, the leading-fermionic retained off-diagonal estimate gives
\[
\eta_{hs}\equiv\left|\frac{\Pi_{hs}^{\mathrm{ret}}(m_h^2)}{a-b}\right|
\simeq1.695\times10^{-5},
\qquad
a=m_h^2-im_h\Gamma_h,\quad b=m_s^2-im_s\Gamma_s.
\]
In the isolated-pole, weak-mixing and SM-dominated-channel limit, the mixing-only leading rate diagnostic is
\[
\delta\mu_h^{(\mathrm{mix,LO})}
\equiv\mu_h^{(\mathrm{mix,LO})}-1
\simeq-\eta_{hs}^2
\simeq-2.87\times10^{-10}.
\]
The corresponding mixing-only fractional rate shift is about \(2.9\times10^{-8}\%\). Here \(\eta_{hs}\) is a first-order off-diagonal mass-insertion diagnostic, not an independently calculated physical mixing angle. The rate estimate is mixing-only: it is not a complete pole signal strength and does not include process-dependent direct-singlet-Yukawa corrections in the \(\mu\) and \(\tau\) channels or the separate searches for the \(100\,\mathrm{GeV}\) scalar.

For comparison, the PDG 2025 update quotes the inclusive global signal strengths \(\mu_{\mathrm{ATLAS}}=1.023^{+0.056}_{-0.053}\) and \(\mu_{\mathrm{CMS}}=1.014^{+0.055}_{-0.053}\) \cite{ParticleDataGroup:2024cfk}. Their uncertainties are approximately \(5.3\%\)--\(5.6\%\), more than \(1.8\times10^8\) times larger than the illustrative \(2.9\times10^{-10}\) mixing-only rate shift. This is a comparison of scales for the isolated \(125\,\mathrm{GeV}\) pole in the stated non-degenerate benchmark, not a model-specific likelihood fit or a bound on the separate \(S\)-like state. The quoted experimental combinations assume a single narrow resonance and fixed signal shapes; they are not a two-scalar fit. The present comparison is deliberately restricted to a point with a well-separated \(S\)-like pole and with \(h\to ss\) closed. For \(m_s<m_h/2\), scalar-pair decays such as \(h\to ss\), when open, depend on the corresponding trilinear coupling and cannot be inferred from \(\cos^2\theta\) alone.

The finite-momentum result does not imply uniformly negligible mixing. The relevant complex inverse propagator is
\[
\mathcal D_B^{-1}(p^2)=
\begin{pmatrix}
p^2-m_h^2+i m_h\Gamma_h&-\Pi_{hs}^{B}(p^2)\\
-\Pi_{sh}^{B}(p^2)&p^2-m_s^2+i m_s\Gamma_s
\end{pmatrix},
\qquad
\det\mathcal D_B^{-1}(p^2)=0.
\]
The near-pole solver retains the momentum dependence of the light-lepton soft \(B_0\) functions and analytically continues them to the resonance sheet. Its hard part uses the source-aligned finite-derivative approximation; at \(p^2=m_h^2\), replacing this by the exact mass-eigenstate hard difference changes the hard term by \(6.11\times10^{-4}\) relative and the total retained \(\Pi_{hs}\) by about \(3.43\times10^{-5}\) relative. The diagonal real self-energies are absorbed into the benchmark mass inputs and the diagonal widths are held fixed, with \(\Gamma_h=4.07\,\mathrm{MeV}\). The resulting poles are therefore first-pass complex-pole diagnostics, not a complete pole-scheme calculation or an experimental lineshape fit.

At exact diagonal-mass degeneracy \(m_s=m_h=125.25\,\mathrm{GeV}\), the two benchmarks give different pole behavior:
\[
\resizebox{\textwidth}{!}{$
\begin{array}{c|cc}
&\text{Benchmark A: }A'A'\text{ open}
&\text{Benchmark B: }A'A'\text{ closed}\\ \hline
\Gamma_s&3.7468\,\mathrm{GeV}&0.15894\,\mathrm{MeV}\\
|\Pi_{hs}/(a-b)|&9.14\times10^{-4}&0.19685\\
(m_{\rm pole},\Gamma_{\rm pole})_{\rm narrow}&
(125.250001\,\mathrm{GeV},4.0727\,\mathrm{MeV})&
(125.250041\,\mathrm{GeV},0.2921\,\mathrm{MeV})\\
(m_{\rm pole},\Gamma_{\rm pole})_{\rm broad}&
(125.264006\,\mathrm{GeV},3746.39\,\mathrm{MeV})&
(125.249959\,\mathrm{GeV},3.9368\,\mathrm{MeV})
\end{array}
$}
\]
For the closed-channel benchmark, a \(0.1\,\mathrm{MeV}\)-step scan over \(m_s\in[125.20,125.30]\,\mathrm{GeV}\) finds a maximum Higgs-like pole-eigenvector composition diagnostic \(|s/h|\simeq0.2039\) near \(m_s=125.2499\,\mathrm{GeV}\); the sampled interval with \(|s/h|\geq0.1\) is approximately \(125.2467\)--\(125.2532\,\mathrm{GeV}\). In the open-channel benchmark the maximum over the same local scan is \(9.15\times10^{-4}\), and no point reaches \(0.1\). The quantity \(|s/h|\) is a non-Hermitian pole-eigenvector diagnostic, not a directly observable mixing angle or rate.

Separately, the exact mass-eigenstate fixed-self-energy scan over the local interval \(m_s\in[10,150]\,\mathrm{GeV}\), with \(1\,\mathrm{MeV}\) spacing and root-refined threshold crossings, gives no interval satisfying \(2|\Pi_{hs}/(a-b)|\geq0.1\) in Benchmark A. In Benchmark B the corresponding diagnostic interval is
\[
m_s\in[125.242553544,125.257446018]\,\mathrm{GeV},
\]
or approximately \(|m_s-m_h|\lesssim7.45\,\mathrm{MeV}\). This is a fixed-self-energy diagnostic, not an allowed or excluded interval. The \(|s/h|\) eigenvector criterion and \(2|\Pi_{hs}/(a-b)|\) criterion are distinct diagnostics and are not identified with one another. In the narrow near-degenerate region, the single-pole approximation \(\mu_h\simeq\cos^2\theta\) is not used; a physical rate requires both pole residues, production and decay amplitudes, interference and the experimental mass response.

The perturbativity condition \(\lambda_S=m_s^2/(2v_s^2)\leq4\pi\) gives
\[
m_s\leq m_s^{\max}=\sqrt{8\pi}\,v_s
=
\begin{cases}
355.4\,\mathrm{GeV},&\text{Benchmark A},\\
1579.2\,\mathrm{GeV},&\text{Benchmark B}.
\end{cases}
\]
The numerical interval \(10\)--\(150\,\mathrm{GeV}\) is a local diagnostic window around the \(125\,\mathrm{GeV}\) Higgs pole, not an exhaustive scan of these full perturbative intervals. For \(m_s\gg m_h\), at fixed benchmark inputs the retained off-diagonal numerator varies slowly with \(m_s\), whereas the diagonal mass separation grows as \(m_s^2\); hence the mixing diagnostic decreases parametrically as \(|\Pi_{hs}|/m_s^2\). The local scan is not used to infer a global scalar-sector exclusion or viability statement. The radial mode \(s\) remains a physical degree of freedom in the UV theory; the scalar-suppressed benchmark only specifies when the axial-current observables in Fig.~2 can be discussed separately from scalar-mediated observables. It does not establish formal decoupling or the absence of scalar effects in every process.

For \(q^2\ll m_s^2,m_{A'}^2\), integrating out \(s\) and comparing local exchange coefficients gives scalar four-lepton operators with coefficient scale
\[
C_{\ell\ell'}^{(s)}\sim\frac{y_{s\ell\ell}y_{s\ell'\ell'}}{m_s^2},
\qquad y_{s\ell\ell}\simeq\frac{m_\ell}{v_s}.
\]
Using \(m_{A'}=2g_Xv_s\), the corresponding \(A'\)-exchange coefficient scale is \(C^{(A')}\sim g_X^2/m_{A'}^2=1/(4v_s^2)\), so
\[
\frac{|C_{\ell\ell'}^{(s)}|}{|C^{(A')}|}
\sim\frac{4m_\ell m_{\ell'}}{m_s^2}.
\]
For \(m_s\geq10\,\mathrm{GeV}\), this coefficient-level ratio is at most \(4.5\times10^{-4}\) for a \(\mu\mu\) operator; at \(m_s=100\,\mathrm{GeV}\) it is \(4.5\times10^{-6}\). For a \(\tau\tau\) operator it can instead be about \(0.13\) at \(m_s=10\,\mathrm{GeV}\). This scaling comparison therefore supports separating scalar exchange from the muon-dominated effective-current benchmark, but it does not justify neglecting scalar effects in every \(\tau\)-sensitive process and is not a process-specific rate comparison.

The one-singlet scalar-suppressed benchmark specifies the regime in which the \(A'\) effective-current observables in Fig.~2 are presented separately from scalar-mediated observables; it does not remove \(s\) from the UV theory, establish formal decoupling, or imply that scalar effects vanish in every process.



\medskip
\noindent {\bf Major Points 2)}. \begin{quote}\itshape
Motivating the range of effective theory parameters. To further aid in the comparison between the UV theory and the two-parameter low-energy theory, it would be useful to explicitly state in one place (perhaps in a table or a figure like Fig. 2) the range of the effective model parameters that are feasible. The main phenomenological results are in Fig. 2, but it was not clear why this range of couplings and \(A'\) masses were chosen: are values outside this range not allowed? What, if anything, does the perturbativity of the UV theory impose on the effective parameters?\end{quote}

\noindent {\bf Reply.}
Fig.~2 displays the intermediate-mass laboratory comparison in the window
\[
10^{-3}\,\mathrm{GeV}\leq m_{A'}\leq10^2\,\mathrm{GeV},
\qquad
10^{-6}\leq g_X\leq1.
\]
This is a plotting window chosen to show complementary experimental coverage, not a universal allowed region or a boundary of the UV model. The motivation for its endpoints, the experimental coverage, and the point-dependent UV consistency conditions are specified below.

\medskip
\noindent {\bf Laboratory coverage within the displayed window.}

Within this window, PIENU probes \(m_{A'}<m_{\pi^+}-m_\mu\simeq33.9\,\mathrm{MeV}\)~\cite{PIENU:2021clt,Altmannshofer:2026opc}, and the indicative NA62 kaon-decay recast extends to \(m_{A'}<m_{K^+}-m_\mu\simeq388\,\mathrm{MeV}\)~\cite{NA62:2021bji}. CCFR and CHARM~II provide complementary neutrino-trident constraints~\cite{CCFR:1991lpl,geiregat1990first}, while NA64\(_\mu\) probes the sub-GeV region~\cite{NA64:2024nwj}. BaBar and Belle~II cover visible resonances up to approximately \(10\) and \(9\,\mathrm{GeV}\), respectively, and the CMS analysis covers \(5\)--\(70\,\mathrm{GeV}\)~\cite{BaBar:2016sci,Belle-II:2024wtd,CMS:2018yxg}. This complementarity, rather than a universal allowed mass interval, motivates the central region displayed.

\medskip
\noindent {\bf Lower-mass region.}

Below the displayed mass window, astrophysical environments provide important additional sensitivity to weakly coupled light mediators. Supernova cooling and energy transport are particularly relevant to muon- and neutrino-coupled bosons~\cite{Blinov:2025aha,Wang:2026supernova}. Their sensitivity also overlaps the sub-GeV part of Fig.~2, so the lower plot edge is not a boundary at which astrophysical constraints suddenly turn on. As explained under Major Point~4.4, the full vector \(L_\mu-L_\tau\) exclusion cannot be copied into the axial model because charged-muon production and absorption differ. These astrophysical probes motivate a separate treatment of the lower-mass region rather than extending the laboratory summary plot over many additional decades.

The more direct reason not to extend the figure to much smaller masses is that the \((g-2)_\mu\) bound already dominates the laboratory constraints displayed at its low-mass end, and its low-mass continuation is fixed by a simple asymptotic behavior. In the scalar-suppressed benchmark, for \(m_{A'}\ll m_\mu\),
\[
\Delta a_\mu^{A'}\simeq-\frac{g_X^2}{4\pi^2}\frac{m_\mu^2}{m_{A'}^2},
\qquad
g_X^{\rm lim}\simeq1.76\times10^{-6}\left(\frac{m_{A'}}{1\,\mathrm{MeV}}\right),
\]
where the second expression uses the paper's approximate \(2\sigma\) requirement \(\Delta a_\mu^{A'}\gtrsim-8.8\times10^{-10}\), obtained from \(\Delta a_\mu=38(63)\times10^{-11}\)~\cite{Leveille:1977rc,Muong-2:2025xyk,Aliberti:2025beg}. The boundary therefore continues towards smaller couplings in proportion to the mediator mass, appearing as a straight line in the log--log plane. Extending Fig.~2 further to the left would chiefly prolong this already dominant precision-bound regime, rather than reveal an additional transition among the laboratory probes compared here. We therefore keep the figure focused on the interval that displays their complementary reach. As a numerical illustration, the continuation crosses the plotted coupling floor \(g_X=10^{-6}\) at approximately \(0.57\,\mathrm{MeV}\). This reasoning assumes the heavy-scalar benchmark, without a cancellation of the negative axial contribution, and does not establish a universal ordering relative to astrophysical bounds.

\medskip
\noindent {\bf Upper-mass region.}

Above the displayed window, the on-shell meson-decay, B-factory, and \(Z\)-decay channels used here have exhausted their kinematic reach. For low-energy virtual exchange, the relevant interaction strength decreases as \(g_X^2/m_{A'}^2\) at fixed coupling, and the heavy-mediator contribution to \((g-2)_\mu\) also decouples as \(m_\mu^2/m_{A'}^2\). Thus these measurements have progressively weaker sensitivity to higher masses at the couplings of interest, which motivates ending this comparison at \(100\,\mathrm{GeV}\). Heavier mediators remain possible targets for higher-energy experiments, including the future muon-collider tests discussed in the paper; they are not excluded or intrinsically unobservable because they lie outside Fig.~2.

\medskip
\noindent {\bf Coupling range and UV perturbativity.}

The coupling endpoints \(10^{-6}\) and \(1\) are likewise display choices: the lower edge retains the low-mass \((g-2)_\mu\) boundary, while the upper edge shows the larger-coupling exclusions. Neither is a strict physical boundary. A conventional gauge perturbativity check is \(|g_X Q_i|\lesssim\sqrt{4\pi}\) for every charged field, evaluated at the scale in question. For the unit-charge leptons this gives \(g_X\lesssim\sqrt{4\pi}\); including the charge-two scalar gives the stronger estimate \(g_X\lesssim\sqrt{4\pi}/2\simeq1.77\). The plotted upper edge lies within both estimates, but neither criterion selects \(g_X=1\) as an endpoint. Perturbativity of the independent UV Yukawa and scalar couplings must also be checked for a chosen completion; it does not follow from the gauge coupling alone. At leading order, the VLL Yukawa parameters are related to the mixing angles by \(\theta_{\ell L}\simeq y_\ell v_h/(\sqrt2M_E)\) and \(\theta_{\ell R}\simeq\lambda_\ell v_s/(\sqrt2M_E)\), as derived under Major Point~1.3. Their perturbativity therefore depends on the chosen \(M_E\), \(v_s\), and mixing-angle inputs; it does not define a universal boundary in the two-parameter \((m_{A'},g_X)\) plane. The explicit \(|\lambda_\ell|\leq\sqrt{4\pi}\) check under Major Point~1.3 applies to its two illustrative portal benchmarks only.

For clarity, the scan variables and their UV interpretation are summarized here:
\begin{center}
\small
\begin{tabular}{@{}P{0.19\linewidth}P{0.34\linewidth}P{0.39\linewidth}@{}}
\hline
Quantity & Range or relation & Interpretation\\
\hline
\(m_{A'}\) & \(10^{-3}\)--\(10^2\,\mathrm{GeV}\) & Main laboratory comparison window; lower masses emphasize astrophysical/precision probes, while the searches shown lose reach at higher masses.\\
\(g_X\) & \(10^{-6}\)--\(1\) & Display window containing the low-mass precision boundary and the larger-coupling exclusions; not a universal allowed interval.\\
\(v_s\) & \(v_s=m_{A'}/(2g_X)\) & Derived from each selected \((m_{A'},g_X)\) point, not an independent scan variable.\\
\(\lambda_S\) & \(\lambda_S=m_s^2/(2v_s^2)=2g_X^2m_s^2/m_{A'}^2\) & Scalar self-coupling in the stated small-portal benchmark; its perturbativity condition must be checked separately.\\
\(m_s\) benchmark & \(m_s\geq10\,\mathrm{GeV}\) for the scalar-suppressed reference case & Reference scalar-mass benchmark used to evaluate the UV consistency condition.\\
\hline
\end{tabular}
\end{center}

In the scalar-mass convention used in the manuscript, imposing \(\lambda_S\leq\lambda_S^{\max}\) gives
\[
m_s\leq\sqrt{2\lambda_S^{\max}}\,v_s.
\]
For the illustrative choice \(\lambda_S^{\max}=4\pi\), this becomes \(m_s\leq\sqrt{8\pi}\,v_s\). Consequently, the reference benchmark \(m_s\geq10\,\mathrm{GeV}\) requires
\[
v_s\geq\frac{10\,\mathrm{GeV}}{\sqrt{8\pi}}
\simeq1.995\,\mathrm{GeV},
\qquad
m_{A'}\gtrsim(3.99\,\mathrm{GeV})\,g_X.
\]
This is the same \(v_s\simeq2\,\mathrm{GeV}\) boundary shown by the dashed line in Fig.~2: the side with \(v_s<2\,\mathrm{GeV}\) does not satisfy the chosen \(m_s\geq10\,\mathrm{GeV}\), \(\lambda_S\leq4\pi\) benchmark. This is a conditional theoretical consistency requirement, not an experimental exclusion. A point that fails it does not support this scalar benchmark; it is not thereby excluded from every possible realization of the UV model if the scalar mass or scalar-sector assumptions are changed. Within the displayed range \(g_X\leq1\), the above gauge-coupling estimates do not set the plot boundary. The scalar self-coupling condition instead imposes a point-dependent check once \(m_s\) is specified.


\medskip
\noindent {\bf Interpretation of Fig.~2.}

The scalar bounds discussed under Major Point~1 constrain \(y_{s\mu\mu}\) directly. They can be translated into the \((m_{A'},g_X)\) plane only for a specified scalar-mass benchmark and under the UV relation \(y_{s\mu\mu}=2g_Xm_\mu/m_{A'}\); they are not direct, benchmark-independent limits on the pure-axial \(A'\) current.

For clarity, Fig.~2 should be read according to the following distinctions:
\begin{center}
\small
\begin{tabular}{p{0.30\linewidth}p{0.62\linewidth}}
\hline
Region or condition & Interpretation \\
\hline
Excluded by a contour shown in Fig.~2 & Excluded under the assumptions of that particular effective-current constraint. \\
Not excluded by the contours shown & Survives the constraints displayed in Fig.~2; this alone does not establish every UV consistency condition. \\
Fails the chosen scalar benchmark & Does not satisfy \(m_s\geq10\,\mathrm{GeV}\) together with \(\lambda_S\leq4\pi\) at that point; other scalar-sector choices are not ruled out by this condition alone. \\
Outside the plotted window & Not forbidden by the scan endpoints themselves; no exclusion or viability claim follows solely from being outside the displayed axes. \\
\hline
\end{tabular}
\end{center}
The contours therefore communicate the constraints actually applied, while the axis endpoints specify the range chosen to present the phenomenology. The table uses the one-VEV matching relation of the original paper. Its two plotting variables parameterize the original \(A'\) benchmark; the table is not a scan over additional UV field content or an enlarged completion parameter space.


\medskip
\noindent {\bf Major Points 3)}.

\begin{quote}\itshape
Are neutrino oscillations possible? The Standard Model is defined with only left-handed neutrinos because of the possibility that neutrinos are their own antiparticle. In such a case neutrino masses are generated through, e.g. through the Weinberg operator.
\end{quote}

\begin{quote}\itshape
Thus it is not obvious to me how neutrino oscillations can be included in this framework. If there is a simple solution, the authors should include it. If there is a known less-simple solution in the literature, the authors should cite this. Otherwise, at best this will be a confusing point. At worst, it would make the construction inconsistent with neutrino oscillation experiments.
\end{quote}

\noindent {\bf Reply.}
The original one-VEV model, with \(S\) as its only \(U(1)_X\)-breaking scalar, does not reproduce the observed three-flavor PMNS mixing: the residual \(Z_2\) separates the electron flavor from the \((\mu,\tau)\) sector. A known singlet-scalar type-I seesaw completion can remove this restriction by adding a charge-one scalar \(\sigma\) and three right-handed neutrinos \(N_R\) to the original field content~\cite{Heeck:2011wj,Asai:2017ryy,Borah:2021twoSingletSeesaw,He:2024twoSingletHNL}. This literature construction applies because \(L_\alpha\) and \(H\) have the same \(U(1)_X\) charges as in the conventional vector \(L_\mu-L_\tau\) realization; the reversed right-handed charged-lepton charges do not enter the neutrino mass operators. The selection rules and explicit mass-generation construction are given in Replies~3.1 and~3.2. The optional completion is subject to the seesaw, perturbativity, and flavor assumptions stated below; it is confined to Major Points~3 and~4.1 and does not change the original one-VEV \(S+\mathrm{VLL}\) model used for non-neutrino phenomenology. The original one-VEV contours are not bounds on the enlarged field content without rematching its couplings and rates.

\medskip
\noindent {\bf Reply to Point 3.1}

\begin{quote}\itshape
In this model, \(U_A(1)_{L_{\mu-\tau}}\) invariance allows Weinberg operators for each generation. However, it appears that both the Yukawa matrices and neutrino masses must be aligned and therefore a non-trivial PMNS matrix does not appear to be possible.
\end{quote}

\noindent {\bf Reply.}
The charge-two scalar alone leaves a restriction on mixing, but the standard charge-one breaking field removes it. Set \(Q_\ell=1\), absorbing the overall charge normalization into \(g_X\). The original assignments are
\[
\begin{gathered}
Q_X(L_e,L_\mu,L_\tau)=(0,1,-1),\qquad
Q_X(e_R,\mu_R,\tau_R)=(0,-1,1),\\
Q_X(H)=0,\qquad Q_X(S)=2.
\end{gathered}
\]
If only \(S\) obtains a new-symmetry-breaking VEV, the phase transformation \(\alpha=\pi\) leaves this VEV invariant and acts as \((+1,-1,-1)\) on the three lepton families. An exact residual \(Z_2\) therefore separates the electron family from the other two.

Define \(\chi_\alpha=L_\alpha\cdot H\equiv\epsilon_{ab}L_\alpha^aH^b\), with the Lorentz-spinor contraction in \(\chi_\alpha\chi_\beta\) understood. Since \(Y(L)=-1/2\) and \(Y(H)=+1/2\), this is the electroweak-invariant Weinberg structure in two-component notation. Its \(U(1)_X\) charge is \(q_\alpha+q_\beta\). The compensating scalar insertions are
\begin{center}
\small
\begin{tabular}{c|c|c|c}
Flavor entry & \(q_\alpha+q_\beta\) & Compensating field & Operator dimension\\
\hline
\(ee,\mu\tau\) & \(0\) & None & 5\\
\(e\mu\) & \(+1\) & \(\sigma^\dagger\) & 6\\
\(e\tau\) & \(-1\) & \(\sigma\) & 6\\
\(\mu\mu\) & \(+2\) & \(S^\dagger\) & 6\\
\(\tau\tau\) & \(-2\) & \(S\) & 6
\end{tabular}
\end{center}
The two entries involving \(\sigma\) are absent in the original field content. No integer number of \(S\) or \(S^\dagger\) insertions compensates their odd charges. Under this \(S\)-only symmetry-breaking assumption, the symmetry-allowed EFT mass matrices have the forms
\[
M_\nu^{(S)}=\begin{pmatrix}a&0&0\\0&b&c\\0&c&d\end{pmatrix},
\qquad
Y_\ell^{(S)}=\begin{pmatrix}y_e&0&0\\0&y_{\mu\mu}&y_{\mu\tau}\\0&y_{\tau\mu}&y_{\tau\tau}\end{pmatrix}.
\]
The off-diagonal \(\mu\)--\(\tau\) charged-lepton Yukawa operators are gauge allowed; setting them to zero is the flavor-diagonal benchmark, not a gauge selection rule. In either case the left-handed rotations retain the electron block and yield \(U_{\rm PMNS}=\operatorname{diag}(1,U_{23})\), up to phases and mass ordering. This need not be the identity, but it cannot describe the observed electron-flavor admixtures. Allowed EFT coefficients alone also do not specify a UV origin of neutrino masses.

Now introduce a complex SM-singlet scalar \(\sigma\) with \(Q_X(\sigma)=1\) and \(\langle\sigma\rangle\ne0\). Its VEV breaks the residual \(Z_2\), and all six independent Majorana entries become allowed. A gauge-invariant EFT parameterization through dimension six is
\[
\mathcal L_{\nu,\mathrm{EFT}}\supset
-\frac12\sum_{\alpha,\beta}K_{\alpha\beta}\chi_\alpha\chi_\beta+\mathrm{h.c.},
\]
\[
K=\begin{pmatrix}
c_{ee}/\Lambda&c_{e\mu}\sigma^\dagger/\Lambda^2&c_{e\tau}\sigma/\Lambda^2\\
c_{e\mu}\sigma^\dagger/\Lambda^2&c_{\mu\mu}S^\dagger/\Lambda^2&c_{\mu\tau}/\Lambda\\
c_{e\tau}\sigma/\Lambda^2&c_{\mu\tau}/\Lambda&c_{\tau\tau}S/\Lambda^2
\end{pmatrix}.
\]
This displays the allowed operators, without assuming independent Wilson coefficients in a particular UV completion or a universal ordering of the heavy and symmetry-breaking thresholds. The seesaw below supplies an explicit renormalizable origin. After the extra \(U(1)\) is broken and the heavy neutrinos are integrated out, the low-energy theory contains the usual Weinberg operator and light Majorana neutrinos. No light right-handed neutrino is required; heavy right-handed fields in the UV are fully compatible with Majorana light neutrinos.

\medskip
\noindent {\bf Reply to Point 3.2}

\begin{quote}\itshape
Perhaps a full model would require right-handed neutrinos. These would now be charged under \(U_A(1)_{L_{\mu-\tau}}\), which likely means the second and third generation must appear as vectorlike combinations. However, it seems to me that mixing between generations is prohibited by \(U_A(1)_{L_{\mu-\tau}}\).
\end{quote}

\noindent {\bf Reply.}
Introduce three SM-singlet right-handed neutrinos with
\[
Q_X(N_{Re},N_{R\mu},N_{R\tau})=(0,1,-1).
\]
These charges allow \(\overline L_\alpha\widetilde H N_{R\alpha}\), with \(\widetilde H=i\sigma_2H^*\). At the renormalizable level with a single neutral electroweak Higgs doublet the Dirac matrix is diagonal. We denote the corresponding neutrino Dirac Yukawa couplings by \(y^\nu_\alpha\), to distinguish them from the charged-lepton Yukawa parameters introduced above. The neutral Majorana bilinears admit bare masses \(M_{ee}\) and \(M_{\mu\tau}\); the two charged singlets supply the other four independent entries. In four-component right-handed-field notation,
\begin{align*}
-\mathcal L_\nu={}&\sum_{\alpha=e,\mu,\tau}y^\nu_\alpha\overline L_\alpha\widetilde H N_{R\alpha}\\
&+\frac12 M_{ee}\overline{N_{Re}^{c}}N_{Re}
+M_{\mu\tau}\overline{N_{R\mu}^{c}}N_{R\tau}\\
&+h_{e\mu}\sigma^\dagger\overline{N_{Re}^{c}}N_{R\mu}
+h_{e\tau}\sigma\overline{N_{Re}^{c}}N_{R\tau}\\
&+\frac12h_{\mu\mu}S^\dagger\overline{N_{R\mu}^{c}}N_{R\mu}
+\frac12h_{\tau\tau}S\overline{N_{R\tau}^{c}}N_{R\tau}+\mathrm{h.c.}
\end{align*}
For example, the total charges of \(\sigma^\dagger\overline{N_{Re}^{c}}N_{R\mu}\) and \(S^\dagger\overline{N_{R\mu}^{c}}N_{R\mu}\) are \(-1+0+1=0\) and \(-2+1+1=0\). All remaining terms are neutral as well. Generation mixing is therefore generated by spontaneous breaking without explicit gauge-symmetry violation.

Take \(\langle H^0\rangle=v_h/\sqrt2\), \(\langle\sigma\rangle=v_\sigma/\sqrt2\), and \(\langle S\rangle=v_s/\sqrt2\), absorbing VEV phases into couplings. Then
\begin{equation}\label{eq:reply-neutrino-matrices}
\begin{gathered}
m_D=\frac{v_h}{\sqrt2}\operatorname{diag}(y^\nu_e,y^\nu_\mu,y^\nu_\tau),\\[3pt]
M_R=\begin{pmatrix}
M_{ee}&h_{e\mu}v_\sigma/\sqrt2&h_{e\tau}v_\sigma/\sqrt2\\
h_{e\mu}v_\sigma/\sqrt2&h_{\mu\mu}v_s/\sqrt2&M_{\mu\tau}\\
h_{e\tau}v_\sigma/\sqrt2&M_{\mu\tau}&h_{\tau\tau}v_s/\sqrt2
\end{pmatrix}.
\end{gathered}
\end{equation}
In the seesaw regime \(\|m_DM_R^{-1}\|\ll1\),
\begin{equation}\label{eq:reply-neutrino-seesaw}
M_\nu\simeq-m_DM_R^{-1}m_D^T.
\end{equation}
All six independent complex entries of \(M_R\) are available. For a nonsingular target \(M_\nu\) and nonsingular diagonal \(m_D\), reconstruct
\begin{equation}\label{eq:reply-neutrino-reconstruction}
M_R=-m_D^T M_\nu^{-1}m_D.
\end{equation}
Each reconstructed entry has a gauge-invariant origin in Eq.~\eqref{eq:reply-neutrino-matrices}. Substitution into the seesaw formula returns the target matrix. This is a constructive mass-texture existence argument, rather than parameter counting alone.

Neglecting small light--heavy nonunitarity corrections in the stated seesaw and VLL limits, define the light charged-lepton rotation \(U_L\). Then
\[
\begin{gathered}
U_\nu^T M_\nu U_\nu=\widehat m,\qquad
U_{\rm PMNS}=U_L^\dagger U_\nu,\\
\widehat m=\operatorname{diag}(m_1,m_2,m_3).
\end{gathered}
\]
For a chosen \(U_{\rm PMNS}\), set \(U_\nu=U_LU_{\rm PMNS}\) and \(M_\nu=U_\nu^*\widehat mU_\nu^\dagger\) in Eq.~\eqref{eq:reply-neutrino-reconstruction}. Nonzero light masses with the observed mass-squared differences and mixing can therefore be realized, even for a nondiagonal charged-lepton rotation. The Majorana parameters are independent of the charged-lepton Yukawa parameters, so no alignment is enforced. Perturbativity, heavy-state masses, and flavor constraints still select the physical parameter region; the construction is not a global viability claim for all UV parameters.

The additional fermions preserve anomaly cancellation. In a left-handed Weyl basis their charges are \((0,-1,+1)\), so
\[
\Delta\mathcal A_{X^3}\propto0^3+(-1)^3+(+1)^3=0,\qquad
\Delta\mathcal A_{\mathrm{grav}^2X}\propto0-1+1=0.
\]
They are SM singlets and add no mixed SM-gauge anomalies. The opposite-charge pair is vectorlike under \(U(1)_X\) before breaking, but the scalar-assisted Majorana entries are allowed after breaking. Adding \(N_R\) without \(\sigma\) would still leave the electron block intact.

The correspondence to existing mechanisms is explicit: identify the charge-\((1,2)\) singlets \((\Phi_1,\Phi_2)\) of Refs.~\cite{Borah:2021twoSingletSeesaw,He:2024twoSingletHNL} with \((\sigma,S)\). Their neutrino interactions then have the form above. See Table~I and Eqs.~(4), (7), (8) of Ref.~\cite{Borah:2021twoSingletSeesaw}, and Sec.~II, Eqs.~(1)--(4) of Ref.~\cite{He:2024twoSingletHNL}; our VEV convention fixes the factors of \(\sqrt2\). This equivalence concerns the neutrino mass sector, not vector-model charged-lepton phenomenology. The single charge-one scalar construction of Ref.~\cite{Asai:2017ryy} instead has \((M_R)_{\mu\mu}=(M_R)_{\tau\tau}=0\), and hence two zeros in \(M_\nu^{-1}\) for diagonal \(m_D\). Our existing \(S\) allows both entries, so that special two-zero-minor prediction is not imposed here.

The additional VEV changes the UV parameter relations. With canonical gauge kinetic terms, before the small kinetic-mixing effects considered separately, the scalar kinetic terms give
\begin{equation}\label{eq:reply-neutrino-gauge-mass}
m_{A'}^2=g_X^2(4v_s^2+v_\sigma^2),\qquad
r_\sigma\equiv\frac{v_\sigma^2}{4v_s^2}.
\end{equation}
Thus
\[
v_s=\frac{m_{A'}}{2g_X\sqrt{1+r_\sigma}},\qquad
y_{s\ell\ell}\simeq\frac{m_\ell}{v_s}
=\frac{2g_Xm_\ell}{m_{A'}}\sqrt{1+r_\sigma},
\]
where the Yukawa expression assumes an approximately \(S\)-radial mass eigenstate and small VLL mixing. These relations describe only the optional completion discussed in this neutrino-specific response. They are not used to redefine the original one-VEV model or the scalar, axial-current, kinetic-mixing, finite-vector, and stellar-cooling analyses elsewhere in the response. A small nonzero \(v_\sigma\) removes the exact selection rule, but the mixing magnitudes must still follow from Eq.~\eqref{eq:reply-neutrino-matrices}, not from small \(r_\sigma\) alone.

The scalar potential permits
\begin{align*}
\Delta V={}&-\mu_\sigma^2|\sigma|^2+\lambda_\sigma|\sigma|^4
+\lambda_{H\sigma}|H|^2|\sigma|^2\\
&+\lambda_{S\sigma}|S|^2|\sigma|^2
+(\kappa S^\dagger\sigma^2+\mathrm{h.c.}),
\end{align*}
where \(\kappa\) has mass dimension one. These interactions affect the scalar spectrum and mixing; nonzero \(\kappa\) removes the otherwise accidental relative scalar-phase symmetry. The relation \(m_s^2=2\lambda_Sv_s^2\) is a reference limit requiring negligible additional mass shifts and mixing. Moreover, \(\sigma\overline E_{\mu L}e_R\) and \(\sigma^\dagger\overline E_{\tau L}e_R\) are gauge allowed. Their coefficients must be sufficiently suppressed to retain the flavor-diagonal charged-lepton benchmark; this is a flavor assumption, not a gauge selection rule.

\medskip
\noindent {\bf 4)}.
\begin{quote}\itshape
Is this the full set of experimental constraints? The second half of the manuscript presents a set of experimental bounds on the two-parameter effective \(A'\) model. It is not obvious if this is the full set of constraints on the model. (See also point \#1 above: I would have liked to see a section discussing the bounds on the \(S\) within this model since it cannot be decoupled.) The following experimental bounds seem like they could be relevant. It would be helpful to either include them or justify why they are subdominant to the experiments presented.
\end{quote}

\noindent {\bf Reply.}

The contours in Fig.~2 constrain the specified two-parameter pure-axial effective-current benchmark; they are not presented as an exhaustive global classification of every point in the UV model. We address the additional issues raised here by summarizing the scalar constraints under Major Point~1, estimating the momentum-dependent kinetic mixing, quantifying the leading effects of the finite vector component, and distinguishing the supernova contributions governed by the unchanged neutrino current from those sensitive to the charged-muon current. The conclusions below are stated for the explicit benchmark and mass regimes in which the corresponding comparison applies.

\medskip
\noindent {\bf 4.1)}.
\begin{quote}\itshape
Neutrino oscillation observations. See point \#3 above. It appears that this model predicts a unit PMNS matrix, which is inconsistent with neutrino oscillations.
\end{quote}

\noindent {\bf Reply.}

The referee's statement is correct for the original \(S\)-only symmetry-breaking limit, but the resulting PMNS matrix is not strictly the identity. With no additional source of residual-\(Z_2\) breaking, the neutrino and charged-lepton matrices have the \(e\oplus(\mu,\tau)\) block structure derived under Major Point~3.1, giving
\[
U_{\rm PMNS}=\operatorname{diag}(1,U_{23}),
\]
up to phases and mass ordering. This permits a \(\mu\)--\(\tau\) rotation but no electron-flavor admixture, and therefore cannot reproduce the observed solar and reactor mixing angles. It is a conditional selection-rule statement for the \(S\)-only vacuum, not a complete neutrino-mass prediction of the original charged-lepton construction.

The optional completion discussed under Major Point~3 retains \(S\), adds a complex singlet \(\sigma\) with \(Q_X(\sigma)=+1\), and introduces three right-handed neutrinos with \(Q_X(N_R)=(0,+1,-1)\). A nonzero \(v_\sigma\) breaks the residual \(Z_2\) and permits the previously forbidden \(e\)--\(\mu\) and \(e\)--\(\tau\) entries of \(M_R\). The Dirac and Majorana matrices are given in Eq.~\eqref{eq:reply-neutrino-matrices}; the charged-lepton-basis relation and constructive reconstruction,
\[
U_{\rm PMNS}=U_L^\dagger U_\nu,\qquad
M_R=-m_D^T M_\nu^{-1}m_D,
\]
show how a general three-flavor mass matrix can be accommodated. The full gauge-invariant operator construction, seesaw condition, and flavor assumptions are detailed under Major Points~3.1 and~3.2. The axial charges of the right-handed charged leptons do not obstruct the construction because those fields do not enter the neutrino Dirac or Majorana mass operators.

The anomaly cancellation and the two-VEV gauge-boson mass relation,
\[
m_{A'}^2=g_X^2(4v_s^2+v_\sigma^2),
\]
are established under Major Point~3.2. This relation applies only to the optional neutrino completion; the original one-VEV model and all non-neutrino responses retain \(m_{A'}=2g_Xv_s\). The one-VEV relation is recovered from the optional completion only in the small-\(r_\sigma\) reference regime specified under Major Point~3. Thus the minimal \(S\)-only limit alone does not reproduce the observed PMNS matrix, whereas the optional \(\sigma+3N_R\) completion provides a gauge-invariant, anomaly-consistent construction that can accommodate it without changing the model baseline used elsewhere in this response.

\medskip
\noindent {\bf 4.2)}.
\begin{quote}\itshape
There is no quantitative estimate of the amount of kinetic mixing generated with the ``benchmark'' choices of \(\theta\leq10^{-2}\). The authors make the insightful point that kinetic mixing goes to zero in the axial limit, but because this limit is tied to the size of the mixing (the mass of the vectorlike leptons), for the above benchmark there is non-zero kinetic mixing. It would be useful to have parametric estimates of the amount of kinetic mixing one expects so that readers may compare to experimental bounds. One could invoke parametric expressions in the literature for the relevant loop diagrams.
\end{quote}

\noindent {\bf Reply.}

Finite light--heavy lepton mixing generates a nonzero photon--\(A'\) mixing of order \(g_X\theta_R^2\). With equal right-handed mixing angles, degenerate physical heavy-lepton masses, and zero independent matching coefficient, \(\theta_R=10^{-2}\) gives \(|\epsilon_{\gamma A'}(0)|=1.44\times10^{-6}|g_X|\). At \(m_{A'}=100\,\mathrm{MeV}\) and \(g_X=0.1\), its on-shell value is \(1.49\times10^{-7}\), well below the \(\mathcal O(10^{-4})\) sensitivity of the NA64 electron-beam invisible search~\cite{NA64:2023lightDM}. The derivation and the conditions for this experimental comparison are given below.

\medskip
\noindent {\bf One-loop origin and matching prescription.}

We define \(\epsilon_{\gamma A'}\) in the photon convention, with induced interaction \(\mathcal L\supset e\epsilon_{\gamma A'}A'_\alpha J_{\rm EM}^\alpha\) after diagonalizing the kinetic terms. It is this coefficient, rather than a hypercharge-mixing coefficient before electroweak symmetry breaking, that enters the dark-photon comparison. In the muon block the chiral couplings are \(g_L^\mu=g_X\) and \(g_R^\mu=-g_X\cos2\theta_{\mu R}\). Consequently,
\[
g_V^\mu=\frac{g_L^\mu+g_R^\mu}{2}=g_X\sin^2\theta_{\mu R},
\qquad
g_A^\mu=\frac{g_R^\mu-g_L^\mu}{2}=-g_X\cos^2\theta_{\mu R}.
\]
The tau block has the opposite charge signs. At one loop, the vector--axial two-point function vanishes in vacuum: its parity-odd tensor is proportional to \(\varepsilon^{\alpha\beta\rho\sigma}q_\rho q_\sigma=0\). Thus the parity-even photon kinetic term depends on \(g_V\), not on the unsuppressed axial coupling. The photon vertex remains diagonal and proportional to the identity in each light--heavy block, so a single off-diagonal \(A'\) vertex cannot close this two-point fermion loop.

The heavy threshold can be exhibited in the same mass basis. Writing \(s_\ell=\sin\theta_{\ell R}\), \(c_\ell=\cos\theta_{\ell R}\), the four diagonal vector couplings are
\[
(g_V^\mu,g_V^\tau,g_V^{E_\mu},g_V^{E_\tau})
=g_X(s_\mu^2,-s_\tau^2,c_\mu^2,-c_\tau^2).
\]
For these leptons of electric charge \(-1\), the one-loop photon mixing in the above convention can be written as
\[
\epsilon_{\gamma A'}(q^2)=\epsilon_{\rm ind}
+\frac{e}{2\pi^2}\sum_{f=\mu,\tau,E_\mu,E_\tau}g_V^f
\int_0^1dx\,x(1-x)
\ln\frac{m_f^2-x(1-x)q^2-i0}{\mu_{\rm r}^2}.
\]
Here \(m_f\) are physical mass eigenvalues, \(\mu_{\rm r}\) is a renormalization scale, and \(\epsilon_{\rm ind}\) is the independent local mixing coefficient. Since \(\sum_f g_V^f=0\), the common ultraviolet term and the explicit \(\mu_{\rm r}\) dependence cancel in this sum; that cancellation does not fix \(\epsilon_{\rm ind}\). We impose the matching condition
\[
\boxed{\epsilon_{\rm ind}(\mu_{\rm match})=0}
\]
at the heavy-VLL matching scale. This boundary condition applies only to the independent local coefficient; it does not set the full momentum-dependent physical mixing \(\epsilon_{\gamma A'}(q^2)\) to zero, because finite light-lepton loop contributions remain. For equal right-handed angles the heavy contribution at zero momentum is
\[
\epsilon_{\rm heavy}(0)=\frac{eg_X\cos^2\theta_R}{12\pi^2}
\ln\frac{m_{E_\mu}^2}{m_{E_\tau}^2}.
\]
It vanishes for degenerate physical heavy states; their full loop functions then cancel at any momentum. Equal coupling weights are also required, and equality of the bare mass parameters alone need not ensure exact degeneracy after lepton mixing. We take this paired-heavy-state reference choice together with the matching condition above. These assumptions leave a finite light-lepton contribution because \(m_\mu\ne m_\tau\).

The difference from imposing only the original common bare mass \(M_E\) can also be estimated. In the small-mixing expansion,
\[
m_{E_\ell}^2\simeq M_E^2(1+\theta_{\ell L}^2+\theta_{\ell R}^2),
\qquad
\epsilon_{\rm heavy}(0)\simeq\frac{eg_X}{12\pi^2}
(\theta_{\mu L}^2-\theta_{\tau L}^2),
\]
where the second expression keeps equal right-handed angles and \(\epsilon_{\rm ind}=0\). With \(|\theta_{\ell L}|\leq10^{-2}\), this residual threshold obeys \(|\epsilon_{\rm heavy}(0)|\lesssim2.56\times10^{-7}|g_X|\), to leading order in the mixing angles. Its momentum dependence is negligible at \(100\,\mathrm{MeV}\) for \(M_E>10\,\mathrm{TeV}\). Thus the physical-mass-degenerate reference and the common-bare-mass choice can be compared without assuming that their thresholds are identical.

\medskip
\noindent {\bf Finite light-lepton loop and momentum dependence.}

For the light-lepton loop, we use the momentum-transfer-dependent kinetic-mixing function studied by Hapitas, Tuckler, and Zhang~\cite{Hapitas:2021kineticMixing}, replacing their charged-lepton vector coupling by \(g_X\sin^2\theta_R\) and analytically continuing their spacelike expression with \(Q^2\to-q^2-i0\). The timelike threshold behavior is compared with Wang and Zhang~\cite{Wang:2026supernova}. Under the matching assumptions above, the remaining contribution is
\[
\epsilon_{\gamma A'}(q^2)=
-\frac{e\,g_X\sin^2\theta_R}{2\pi^2}
\int_0^1 dx\,x(1-x)
\log\!\left[
\frac{m_\tau^2-x(1-x)q^2-i0}
     {m_\mu^2-x(1-x)q^2-i0}
\right],
\]
where \(\theta_{\mu R}=\theta_{\tau R}\equiv\theta_R\). The light-lepton ultraviolet terms cancel separately in this limit. At \(q^2=0\), \(\int_0^1dx\,x(1-x)=1/6\), giving
\[
\epsilon_{\gamma A'}(0)=
-\frac{e\,g_X\sin^2\theta_R}{12\pi^2}
\ln\frac{m_\tau^2}{m_\mu^2}
=-\frac{e\,g_X\sin^2\theta_R}{6\pi^2}
\ln\frac{m_\tau}{m_\mu}.
\]
Using \(e=0.3028\), \(m_\mu=105.658~\mathrm{MeV}\), \(m_\tau=1.77686~\mathrm{GeV}\), and \(\theta_R=10^{-2}\), the zero-momentum magnitude is
\[
\begin{aligned}
|\epsilon_{\gamma A'}(0)|&\simeq1.44\times10^{-6}|g_X|\quad(\theta_R=10^{-2}),\\
&\simeq1.44\times10^{-7}
\left(\frac{|g_X|}{0.1}\right)
\left(\frac{\theta_R}{10^{-2}}\right)^2
\quad\text{for small }\theta_R.
\end{aligned}
\]
At fixed off-diagonal lepton mass entry \(m_R\), \(\theta_R\simeq m_R/M_E\), so this contribution scales as \(g_Xm_R^2/M_E^2\). This explicitly exhibits its decoupling with the vector-like-lepton mass under that fixed-parameter choice; maintaining the observed light-lepton masses can require correlated changes of the Yukawa parameters.

For an on-shell \(A'\), the relevant timelike argument is \(q^2=m_{A'}^2\). The loop function becomes complex above the charged-lepton pair thresholds; its magnitude is not constant across the scan. A numerical evaluation of the integral gives the following representative values for \(\theta_R=10^{-2}\):
\begin{center}
\small
\begin{tabular}{c|c|c}
\(m_{A'}\) &
\shortstack{\(\left|\epsilon_{\gamma A'}(m_{A'}^2)\right|\)\\ \(/\left|g_X\right|\)} &
\shortstack{\(\left|\epsilon_{\gamma A'}(m_{A'}^2)\right|\)\\ \(/\left|\epsilon_{\gamma A'}(0)\right|\)}\\
\hline
\(10~\mathrm{MeV}\) & \(1.44\times10^{-6}\) & \(1.00\)\\
\(100~\mathrm{MeV}\) & \(1.49\times10^{-6}\) & \(1.04\)\\
\(200~\mathrm{MeV}\) & \(1.82\times10^{-6}\) & \(1.26\)\\
\(212~\mathrm{MeV}\) & \(2.12\times10^{-6}\) & \(1.47\)\\
\(300~\mathrm{MeV}\) & \(1.74\times10^{-6}\) & \(1.21\)\\
\(500~\mathrm{MeV}\) & \(1.40\times10^{-6}\) & \(0.97\)\\
\(1~\mathrm{GeV}\) & \(1.08\times10^{-6}\) & \(0.75\)\\
\(3~\mathrm{GeV}\) & \(8.06\times10^{-7}\) & \(0.56\)\\
\(10~\mathrm{GeV}\) & \(5.47\times10^{-8}\) & \(3.79\times10^{-2}\)\\
\(100~\mathrm{GeV}\) & \(4.84\times10^{-10}\) & \(3.35\times10^{-4}\)
\end{tabular}
\end{center}
The value near \(m_{A'}=2m_\mu\) agrees with the approximately \(1.47\) enhancement of the magnitude over its zero-momentum value reported by Wang and Zhang~\cite{Wang:2026supernova}. For the low-mass benchmark, the induced electron-current coupling is of order \(e|\epsilon|\simeq4.4\times10^{-7}|g_X|\). The momentum-dependent expression should be used for on-shell production and decay; spacelike virtual exchange requires evaluating the same function at the relevant spacelike \(q^2\).

For the reference normalization \(g_X=0.1\), the zero-momentum value is \(1.44\times10^{-7}\), and the largest tabulated on-shell value is \(2.12\times10^{-7}\). The displayed range in Fig.~2 extends to \(g_X=1\), so the corresponding conservative upper estimates over that rectangle are ten times larger. These estimates do not assert that the largest-coupling points satisfy the other constraints.

\medskip
\noindent {\bf NA64 comparison and decay branching fractions.}

The appropriate dark-photon comparison is the NA64 electron-beam missing-energy search~\cite{NA64:2023lightDM}, rather than its thermal-dark-matter parameter combinations. The experiment constrains the electromagnetic production vertex \(e\epsilon_{\gamma A'}A'_\alpha J_{\rm EM}^\alpha\), which is precisely the vertex induced by the photon--\(A'\) kinetic mixing here. Its Fig.~3 gives a 90\% confidence-level exclusion boundary in the \((m_{A'},\epsilon)\) plane from \(9.37\times10^{11}\) electrons on target. We compare to that mass-dependent boundary, not to the thermal-relic parameters \(y\), \(\alpha_D\), or \(m_\chi\).

In the original model, neglecting neutrino masses, the leading widths below the dimuon threshold are
\[
\Gamma_{\nu\bar\nu}=\frac{g_X^2m_{A'}}{12\pi},
\qquad
\Gamma_{e^+e^-}=\frac{e^2|\epsilon_{\gamma A'}(m_{A'}^2)|^2m_{A'}}{12\pi}
(1+2r_e)\sqrt{1-4r_e},
\]
where the first expression sums the two charged neutrino flavors and \(r_e=m_e^2/m_{A'}^2\); the electron width is zero below threshold. In the original model, this sum runs over the two massless active-flavor channels \(\nu_\mu\) and \(\nu_\tau\), with \(U(1)_X\) charges \(+1\) and \(-1\), respectively; the electron flavor is uncoupled. At \(m_{A'}=100\,\mathrm{MeV}\) and \(\theta_R=10^{-2}\), \(\Gamma_{e^+e^-}/\Gamma_{\nu\bar\nu}\simeq2.05\times10^{-13}\), so \(\operatorname{Br}_{\rm inv}\simeq1\).

Above the dimuon threshold, the direct axial muon current opens a visible channel. In the same convention,
\[
\Gamma_{\mu^+\mu^-}
=\frac{m_{A'}}{12\pi}\sqrt{1-4r_\mu}
\left[(g_V^\mu)^2(1+2r_\mu)+(g_A^\mu)^2(1-4r_\mu)\right],
\qquad r_\mu=\frac{m_\mu^2}{m_{A'}^2}.
\]
At \(m_{A'}=300\,\mathrm{MeV}\), with \(g_V^\mu=g_X\sin^2\theta_R\), \(g_A^\mu=-g_X\cos^2\theta_R\), and \(\theta_R=10^{-2}\), this gives
\[
\frac{\Gamma_{\mu^+\mu^-}}{\Gamma_{\nu\bar\nu}}\simeq0.3576,
\qquad
\operatorname{Br}_{\rm inv}\simeq0.7366.
\]
The loop-induced electromagnetic widths, including the hadronic width above the two-pion threshold, are negligible relative to these direct widths at the stated benchmark. We therefore compare the missing-energy signal using
\[
\epsilon_{\rm eff}(m_{A'})
=|\epsilon_{\gamma A'}(m_{A'}^2)|
\sqrt{\operatorname{Br}_{\rm inv}(m_{A'})},
\qquad
N_{\rm sig}\propto\epsilon_{\rm eff}^2,
\]
assuming the published missing-energy acceptance for an invisible spin-one state. This is a leading-yield rescaling, not a detector-level recast of visible dimuon decays.

For a transparent mass-dependent comparison, we read the NA64 90\% CL boundary from Fig.~3 of Ref.~\cite{NA64:2023lightDM} in log--log coordinates. The figure does not provide a machine-readable table; we therefore quote conservative visual-read-off intervals, with an approximately \(0.2\)--\(0.3\) dex line-thickness/read-off uncertainty, rather than spurious high-precision point values. The lower edge of each interval is used in the ratios below. The model values are evaluated from the full timelike loop function at the same mass:
\begin{center}
\small
\resizebox{\textwidth}{!}{%
\begin{tabular}{c|c|c|c|c|c}
\(m_{A'}\) &
\( |\epsilon|/|g_X| \) &
\(\operatorname{Br}_{\rm inv}\) &
\(\epsilon_{\rm NA64}^{90\%}\) (visual range) &
\(\epsilon_{\rm eff}\) at \(g_X=0.1\) &
\(R\) at \(g_X=0.1\)\\
\hline
\(10~\mathrm{MeV}\) & \(1.44\times10^{-6}\) & \(1.000\) &
\((0.8\text{--}2.0)\times10^{-5}\) & \(1.44\times10^{-7}\) & \(1.80\times10^{-2}\)\\
\(100~\mathrm{MeV}\) & \(1.49\times10^{-6}\) & \(1.000\) &
\((0.5\text{--}1.2)\times10^{-4}\) & \(1.49\times10^{-7}\) & \(2.99\times10^{-3}\)\\
\(200~\mathrm{MeV}\) & \(1.82\times10^{-6}\) & \(1.000\) &
\((0.7\text{--}1.8)\times10^{-4}\) & \(1.82\times10^{-7}\) & \(2.59\times10^{-3}\)\\
\(300~\mathrm{MeV}\) & \(1.74\times10^{-6}\) & \(0.737\) &
\((1.0\text{--}2.5)\times10^{-4}\) & \(1.50\times10^{-7}\) & \(1.50\times10^{-3}\)
\end{tabular}%
}
\end{center}
Here \(R\equiv\epsilon_{\rm eff}/\epsilon_{\rm NA64,low}\), where \(\epsilon_{\rm NA64,low}\) is the conservative lower edge of the visually read-off interval. The corresponding signal-yield ratios \(R^2\) are \(3.3\times10^{-4}\), \(8.9\times10^{-6}\), \(6.7\times10^{-6}\), and \(2.2\times10^{-6}\), respectively, at \(g_X=0.1\). Since \(\epsilon\propto g_X\), the amplitude ratios are ten times larger and the yield ratios one hundred times larger at \(g_X=1\); even there the largest conservative ratio is \(R\simeq0.18\) at \(10\,\mathrm{MeV}\), corresponding to about \(3.3\%\) of the limiting signal yield. The other three mass points remain below \(10^{-3}\) of the limiting yield at \(g_X=1\). These upper-edge values are a conservative comparison over the plotted rectangle, not a claim that those large-coupling points satisfy every other constraint.

The secondary-positron resonance energy is \(E_{e^+}^{\rm res}=m_{A'}^2/(2m_e)\): it is approximately \(0.098\), \(9.78\), and \(39.1\,\mathrm{GeV}\) at 10, 100, and 200 MeV, respectively, below the approximately \(50\,\mathrm{GeV}\) missing-energy selection; at 300 MeV it is approximately \(88.1\,\mathrm{GeV}\), so the published NA64 curve includes the resonant secondary-positron contribution as well as bremsstrahlung. The rescaling by \(\operatorname{Br}_{\rm inv}\) accounts for the open dimuon width at this point. For \(g_X=0.1\) and 1 the total width remains narrow compared with the mass. No dark-matter particle or thermal-relic assumption is required; the invisible final states here are neutrinos.

The zero-momentum light-lepton loop gives \(|\epsilon_{\gamma A'}(0)|=1.44\times10^{-6}|g_X|\), whereas the paired-heavy-mass on-shell light-lepton contribution at \(m_{A'}=100\,\mathrm{MeV}\) is \(1.49\times10^{-6}|g_X|\). The residual common-bare-mass threshold contribution is at most \(2.56\times10^{-7}|g_X|\). Applying the triangle inequality, and using the negligible momentum dependence of this heavy-threshold term at 100 MeV for \(M_E>10\,\mathrm{TeV}\), gives the conservative leading-order on-shell bound
\[
\left|\epsilon_{\gamma A'}(m_{A'}^2)\right|_{m_{A'}=100\,\mathrm{MeV}}
\leq (1.49+0.256)\times10^{-6}|g_X|
\simeq1.75\times10^{-6}|g_X|.
\]
The mass-dependent table uses the paired physical-heavy-mass reference; this residual-threshold estimate does not change the conclusion that the induced electron-production channel is below the published NA64 sensitivity at the representative mass points. This comparison is separate from the direct muon-production constraint of NA64\(_\mu\) in Fig.~2. 

Unequal right-handed mixing angles, independent heavy-mass splittings, or a nonzero independent matching coefficient require evaluating the general four-state expression above.

\medskip
\noindent {\bf 4.3)}.
\begin{quote}\itshape
Another way of reframing the above point: the bounds presented are in the pure axial limit. Since the authors present a UV model that only realizes the axial coupling in a decoupling limit, it would be useful to comment on possible experimental bounds on the small vector-like coupling in a model with non-infinite vectorlike fermion masses. If these bounds are always small compared to the axial bounds, this should be stated and justified. The importance of this point is as follows: often bounds are sensitive to the chiral structure of interactions. Sometimes a vector interaction is more tightly constrained than its axial counterpart. Since the authors focus on the phenomenology in the axial limit, one wants to make sure that any small-but-finite vector interaction isn't significant because the constraints on the coupling are orders-of-magnitude larger than the axial coupling.
\end{quote}

\noindent {\bf Reply.}

Within the adopted finite-mixing domain, the vector component is retained but is parametrically and numerically subleading to the axial component in the observables considered. The conclusion follows from the common-\(g_X\) coupling hierarchy and the process-specific comparisons below; it is not extended to arbitrary mixing angles or arbitrarily close to the dimuon threshold.

\medskip
\noindent {\bf Finite-mixing domain and vector-to-axial hierarchy.}

We restrict the UV parameter region underlying the axial-vector phenomenology to the finite-mixing domain
\[
|\theta_{\ell L}|,\ |\theta_{\ell R}|\leq10^{-2},\qquad \ell=\mu,\tau,
\]
under the small-mixing condition adopted for the UV construction. This is a finite-mass domain, not the strict \(M_E\to\infty\) limit: in the seesaw-like regime,
\[
m_\ell\simeq\theta_{\ell L}\theta_{\ell R}M_{E\ell}.
\]
For the degenerate-VLL reference choice, saturating both angle upper values in the \(\tau\) sector corresponds to \(M_E\simeq17.8~\mathrm{TeV}\); smaller mixing angles correspond to larger, still finite VLL masses. In the mass basis, the diagonal light-lepton couplings retain the opposite overall \(U(1)_X\) signs of the two flavor sectors:
\[
\begin{aligned}
(g_V^\mu,g_A^\mu)&=g_X\bigl(\sin^2\theta_{\mu R},-\cos^2\theta_{\mu R}\bigr),\\
(g_V^\tau,g_A^\tau)&=-g_X\bigl(\sin^2\theta_{\tau R},-\cos^2\theta_{\tau R}\bigr).
\end{aligned}
\]
Their absolute vector-to-axial ratios are nevertheless the same functions of the corresponding right-handed angles,
\[
\left|\frac{g_V^\ell}{g_A^\ell}\right|
=\tan^2\theta_{\ell R}\leq1.01\times10^{-4},
\qquad \ell=\mu,\tau.
\]
Thus the magnitudes of the vector couplings are suppressed by at least four orders of magnitude relative to the corresponding axial couplings throughout the parameter region considered. At the reference value \(g_X=0.1\), this gives \(|g_V^\ell|\leq10^{-5}\); at the plotting upper edge \(g_X=1\), it gives \(|g_V^\ell|\leq10^{-4}\). The neutrino couplings remain \(g_{\nu_\mu}=+g_X\) and \(g_{\nu_\tau}=-g_X\).

\medskip
\noindent {\bf Criterion for vector-specific experimental bounds.}

The angle restriction also gives a direct criterion for assessing a vector-specific experimental bound. At a fixed \(m_{A'}\), let \(g_{V,\mathrm{lim}}\) and \(g_{A,\mathrm{lim}}\) denote upper limits on the magnitudes of the corresponding vector and axial charged-lepton couplings, obtained with matched production, decay, and acceptance assumptions. Since both couplings are proportional to the same \(g_X\), the vector limit is weaker than the axial limit when
\[
\frac{g_{V,\mathrm{lim}}}{g_{A,\mathrm{lim}}}>
\tan^2\theta_{\ell R}.
\]
Equivalently, the angle required for the vector-specific bound to remain subdominant is
\[
|\theta_{\ell R}|<
\arctan\sqrt{\frac{g_{V,\mathrm{lim}}}{g_{A,\mathrm{lim}}}}.
\]
For our adopted upper value \(|\theta_{\ell R}|=10^{-2}\), a vector-specific limit would have to constrain the vector coupling to below approximately \(10^{-4}\) of the corresponding axial limit—more than four orders of magnitude tighter at the same mass—to yield a stronger bound on \(g_X\). This is the quantitative criterion defining the small-mixing region used here; it does not assert that vector-sensitive limits vanish for arbitrary mixing angles or for unmatched experimental assumptions.

For neutrino trident production, the ratio of the leading Standard-Model interference terms induced by the vector and axial components is parametrically
\[
\left|\frac{C_V\delta C_V}{C_A\delta C_A}\right|
\sim\frac{C_V}{C_A}\tan^2\theta_{\mu R}
\lesssim2.0\times10^{-4},
\qquad C_V\simeq0.96,\quad C_A=\frac12,
\]
at the same momentum transfer, before the distinct finite-\(m_\mu\) response functions are included. For the muon anomalous magnetic moment,
\[
\left|\frac{\Delta a_\mu^{V}}{\Delta a_\mu^{A}}\right|
=\tan^4\theta_{\mu R}\frac{F_V(r)}{F_A(r)}
\lesssim2.1\times10^{-9},
\]
using \(F_V/F_A\leq0.20\) over the mass range considered. These are ratios of contributions to the specified observables, not standalone experimental exclusion limits; they show that the finite vector component is subleading in the trident-interference and \((g-2)_\mu\) observables used in our analysis.

\medskip
\noindent {\bf Comparison with the four-muon searches.}

For the BaBar, Belle~II, and CMS four-muon final-state-radiation searches used in Fig.~2, the unit-coupling vector and axial production kernels become nearly identical for \(m_{A'}\gg m_\mu\), with correspondingly similar acceptances; away from threshold, the dimuon widths and branching fractions also approach the same leading behavior~\cite{BaBar:2016sci,Belle-II:2024wtd,CMS:2018yxg}. Thus the vector-only coupling reach in these same channels is of the same order as the axial reach, rather than more than four orders of magnitude stronger. In the mixed current, the separate vector-current-squared production contribution is weighted by
\[
\left|\frac{g_V^\mu}{g_A^\mu}\right|^2
=\tan^4\theta_{\mu R}\leq1.01\times10^{-8},
\]
in this away-from-threshold regime. The dimuon branching fraction still requires the threshold-dependent width factors given below.

The loop-induced electromagnetic coupling associated with the same vector admixture is treated separately through the kinetic-mixing estimate in our response to Major Point~4.2.

\medskip
\noindent {\bf Threshold-dependent dimuon widths.}

For visible dimuon decays, the vector and axial partial widths have different threshold behavior:
\[
\Gamma_V\propto(g_V^\mu)^2(1+2r_\mu)\sqrt{1-4r_\mu},\qquad
\Gamma_A\propto(g_A^\mu)^2(1-4r_\mu)^{3/2},
\quad r_\mu=\frac{m_\mu^2}{m_{A'}^2}.
\]
Their ratio is therefore
\[
\frac{\Gamma_V}{\Gamma_A}
=\tan^4\theta_{\mu R}\,
\frac{1+2r_\mu}{1-4r_\mu}.
\]
For any chosen tolerance \(\eta\) on the relative vector contribution, the corresponding angle condition is mass dependent:
\[
|\theta_{\mu R}|<
\arctan\!\left[
\eta^{1/4}
\left(\frac{1-4r_\mu}{1+2r_\mu}\right)^{1/4}
\right].
\]
Away from threshold, the vector-squared contribution is suppressed by \(\tan^4\theta_{\mu R}\leq1.01\times10^{-8}\). Arbitrarily close to \(m_{A'}=2m_\mu\), however, the axial phase-space factor vanishes more rapidly, so no nonzero mass-independent upper bound on the mixing angle can keep the vector-to-axial width ratio uniformly small arbitrarily close to threshold.

Thus, under the adopted small-mixing condition \(|\theta_{\ell R}|\leq10^{-2}\) and away from the dimuon threshold, the process-specific estimates show that the finite vector component is subleading in neutrino trident production and \((g-2)_\mu\), while the vector-only and axial-only coupling reaches in the high-mass four-muon searches are of the same order. Near \(m_{A'}=2m_\mu\), however, the axial partial width is more strongly phase-space suppressed, so the visible dimuon rate must be evaluated with the finite-angle widths rather than a uniform pure-axial approximation.

\noindent {\bf 4.4)}.
\begin{quote}\itshape
Are stellar cooling bounds not relevant? See e.g. 2511.09619 for a recent paper on this. If they are not, then this is certainly worth justifying.
\end{quote}

\noindent {\bf Reply.}

Supernova cooling is relevant to the sub-GeV region. The referee-cited analysis is now published as JHEP 06 (2026) 155~\cite{Blinov:2025aha}; a recent update for the conventional vector \(L_\mu-L_\tau\) boson, using the SFHo-18.8 profile, reports SN1987A sensitivity up to \(m_{A'}\sim300~\mathrm{MeV}\) and low-energy-supernova sensitivity up to about \(600~\mathrm{MeV}\), with SN1987A coupling sensitivity of order \(10^{-10}\) and low-energy-supernova sensitivity reaching \(6.4\times10^{-12}\)~\cite{Wang:2026supernova}. These values provide reference scales for assessing the axial benchmark. Kahn et al. give an earlier discussion of stellar cooling for ultralight axial forces with \(m_{A'}<2m_e\): their Sec.~2.4 emphasizes that an axial \(A'\) can have unsuppressed neutrino couplings and a stellar production mechanism not necessarily dominated by kinetic mixing, and that a dedicated analysis is required~\cite{Kahn:2016vjr}. This motivates treating the axial case separately, but does not provide a model-specific SN1987A contour for the present \(L_\mu-L_\tau\) completion.

\medskip
\noindent {\bf Neutrino-coalescence production and invisible decays.}

In our model, the tree-level neutrino couplings are the same as in the conventional \(L_\mu-L_\tau\) model,
\[
g_{\nu_\mu}=+g_X,\qquad g_{\nu_\tau}=-g_X.
\]
The flavor-summed neutrino-coalescence matrix element and the invisible partial widths \(A'\to\nu_\alpha\bar\nu_\alpha\) therefore have the same coupling normalization as in the conventional vector \(L_\mu-L_\tau\) reference at fixed \((m_{A'},g_X)\), for the same neutrino phase-space distributions and supernova profile. This follows directly from the original model's active-flavor gauge current \(g_{\nu_\mu}=+g_X\), \(g_{\nu_\tau}=-g_X\); no additional neutrino fields or completion assumptions are used in this current-level comparison. In the cited supernova profile, neutrino coalescence dominates the production luminosity for \(m_{A'}\gtrsim10~\mathrm{MeV}\). Accordingly, for \(10~\mathrm{MeV}\lesssim m_{A'}<2m_\mu\), the published SN1987A cooling result is a useful reference scale: neutrino-coalescence production is unchanged and dominates the luminosity in the cited profile, while \(A'\) decays predominantly to neutrinos. We use this as a benchmark comparison, not as a fully recast exclusion contour for the axial model. Equality of the neutrino-coalescence matrix element and invisible partial-width normalization is a current-level statement; it does not imply equality of the complete supernova cooling or energy-transport limits. Those limits also depend on inverse absorption, optical depths and trapping, decay and energy deposition, and the stellar profile; the charged-muon axial and longitudinal channels discussed below can change these parts of the transport problem.

\medskip
\noindent {\bf Charged-muon axial semi-Compton production and absorption.}

The distinction arises in processes involving the charged-muon current. The semi-Compton process \(\mu\gamma\to\mu A'\) and its inverse absorption depend on whether that current is vector or axial. For on-shell muons, the vector current is conserved whereas the axial current obeys
\[
q_\alpha\bar u(p')\gamma^\alpha u(p)=0,\qquad
q_\alpha\bar u(p')\gamma^\alpha\gamma_5u(p)
=2m_\mu\bar u(p')\gamma_5u(p),
\qquad q=p'-p.
\]
The non-conserved axial current therefore contributes to longitudinal \(A'\) emission. In the relativistic longitudinal limit \(E_{A'}\gg m_{A'}\), Goldstone equivalence gives the effective coupling scale
\[
|g_{G\mu}|\sim\frac{2|g_A^\mu|m_\mu}{m_{A'}}.
\]
The semi-Compton emissivity and inverse absorption must consequently be evaluated with the axial current and its polarization structure; the vector-model curve is not directly transferable in the low-mass region where semi-Compton production dominates. Above \(m_{A'}=2m_\mu\), the axial and vector dimuon widths also differ in their threshold dependence, as quantified under Major Point~4.3, affecting visible-decay and low-energy-supernova energy-deposition constraints.

\medskip
\noindent {\bf Kinetic-mixing-induced electromagnetic processes.}

The low-momentum kinetic-mixing estimate has a separate scope. In the equal-mixing benchmark,
\[
\frac{|\epsilon_{\gamma A'}(0)|}{|g_X|}\simeq1.4\times10^{-6},
\]
compared with approximately \(1/70\) in the conventional vector model. At the same vacuum kinematics, this suppresses only electromagnetic-current-mediated amplitudes and rates, such as plasmon conversion or the mixing-induced proton coupling in \(np\) bremsstrahlung, by about \(10^{-4}\) and \(10^{-8}\), respectively. It does not suppress neutrino coalescence or the direct axial semi-Compton process. These factors are same-kinematics vacuum scaling estimates, not supernova emissivity ratios; an actual in-medium rate requires the momentum-dependent mixing and plasma propagators.

\medskip
\noindent {\bf Interpretation and applicability.}

Thus stellar constraints are not dismissed: the neutrino-coalescence contribution is applicable where it dominates, while the full vector-current supernova exclusion curve is not treated as a universal exclusion of the axial model. The semi-Compton and electromagnetic-current-mediated channels are distinguished according to their respective interactions, preserving the relevant astrophysical information without applying a bound outside the interaction structure for which it was derived.

\medskip
\noindent {\bf Minor Point 1)}. \begin{quote}\itshape
Range of parameters. The general construction of the model is clearly written, but it is difficult for the reader to envision the mass and coupling range that the authors are targeting. It would be very helpful if the authors decide on a range of \(m_{A'}\) and \(g_X\) of interest. For example: is there a range of masses where experimental or theoretical limits are less constraining? Is there some set of phenomenology that the authors seek to realize with a ``light'' \(A'\)? If so, they should define ``light.'' (See, e.g. the first word of the introduction, which is not defined with respect to any reference mass scale.)\end{quote}

\noindent {\bf Reply.}
The word "light" is not intended to specify a separate quantitative mass regime. The Abstract describes the broader MeV-to-TeV scale, while Fig.~2 displays the specific comparison window
\[
10^{-3}\,\mathrm{GeV}\leq m_{A'}\leq10^2\,\mathrm{GeV},
\qquad
10^{-6}\leq g_X\leq1.
\]
The numerical range and the observables included in the analysis, rather than the adjective "light", define the phenomenological domain considered here. The lower end retains the low-mass region in which the laboratory \((g-2)_\mu\) constraint is already important; its asymptotic continuation and the overlapping, model-dependent astrophysical sensitivity are discussed under Major Points~2 and~4.4. At the upper end, the meson-decay, B-factory, and \(Z\)-decay searches included in Fig.~2 approach or exhaust their kinematic coverage, while heavier mediators require higher-energy probes. Neither endpoint therefore marks a physical boundary of the model or a universal exclusion boundary. The detailed motivation for the window and the complementary experimental coverage is given in our response to Major Point~2.

\medskip
\noindent {\bf Minor Point 2)}.
\begin{quote}\itshape
Defining the pure-axial limit. Equation (12) gives the coupling of the \(A'\) to the mass-basis fields. It is here that the ``pure axial coupling'' limit is defined? i.e. \(\theta_{\mu R}\to0\)? If that is the case, it would help to define this as the pure-axial limit, a phrase which is used below (36) but not otherwise quantified.
\end{quote}

\noindent {\bf Reply.}

We agree that the pure-axial limit should be quantified when the mass-basis couplings are introduced in Eq.~(12). In the convention used there, the strict pure-axial limit is \(\theta_{\mu R}\to0\):
\[
g_V^\mu=g_X\sin^2\theta_{\mu R}\to0,
\qquad
g_A^\mu=-g_X\cos^2\theta_{\mu R}\to-g_X.
\]
For finite \(\theta_{\mu R}\), the vector coupling is nonzero, so the interaction is not strictly purely axial. In the tau sector, the opposite \(U(1)_X\) charge assignment gives the opposite overall sign, while the coupling magnitudes are \(|g_V^\tau|=|g_X|\sin^2\theta_{\tau R}\) and \(|g_A^\tau|=|g_X|\cos^2\theta_{\tau R}\). Thus \(\left|g_V^\tau/g_A^\tau\right|=\tan^2\theta_{\tau R}\), and the strict pure-axial limit of the full charged-lepton current requires both right-handed mixing angles to vanish. Finite right-handed mixing angles describe a small vector admixture rather than the strict pure-axial limit.

\medskip
\noindent {\bf Minor Point 3)}.
\begin{quote}\itshape
Bounds on the vectorlike leptons. Because the mixing angles scale inversely with the vectorlike lepton masses, it would be helpful to present lower limits on these masses. This, in turn, gives an upper limit on the \(\theta_{(\mu,\tau)R}\) mixing angles which then gives an upper limit on the \(A'\) vectorlike coupling to \(L_{\mu-\tau}\). How do these bounds compare with the \(\theta_{(\mu,\tau)(L,R)}\) bounds in equation (27)?
\end{quote}

\noindent {\bf Reply.}

In the decoupling regime, the light--heavy mixing angles scale inversely with the vector-like-lepton masses at fixed off-diagonal mass parameters. After matching to the observed charged-lepton masses, Eqs.~(21) and (23) give
\[
m_\ell\simeq\theta_{\ell L}\theta_{\ell R}M_{E\ell},
\qquad \ell=\mu,\tau.
\]
Combining these relations with the small-mixing condition adopted in Eq.~(27),
\[
|\theta_{\ell L}|,\;|\theta_{\ell R}|\leq10^{-2},
\]
gives the following conditional mass scales at leading order:
\[
M_{E\mu}\gtrsim\frac{m_\mu}{(10^{-2})^2}\simeq1.06~\mathrm{TeV},
\qquad
M_{E\tau}\gtrsim\frac{m_\tau}{(10^{-2})^2}\simeq17.8~\mathrm{TeV}.
\]
For the degenerate choice \(M_{E\mu}=M_{E\tau}\equiv M_E\) in Eq.~(22), the tau relation therefore gives the stronger conditional requirement \(M_E\gtrsim17.8~\mathrm{TeV}\). This follows from imposing the mixing-angle condition and reproducing the charged-lepton masses simultaneously; it is not an independent direct-search mass limit. The separate condition \(M_E>10~\mathrm{TeV}\) in Eq.~(22) is weaker than this combined requirement.

Equation~(27) and the mass relations play different roles: Eq.~(27) bounds each mixing angle individually in the adopted small-mixing setup, whereas the mass relations fix the products \(\theta_{\ell L}\theta_{\ell R}\). A mass relation by itself would not yield an individual upper bound on \(\theta_{\ell R}\). Using the Eq.~(27) condition in the mass-basis couplings,
\[
\begin{gathered}
|g_V^\mu|=|g_X|\sin^2\theta_{\mu R},\qquad
|g_A^\mu|=|g_X|\cos^2\theta_{\mu R},\\
|g_V^\tau|=|g_X|\sin^2\theta_{\tau R},\qquad
|g_A^\tau|=|g_X|\cos^2\theta_{\tau R},\\
\left|\frac{g_V^\ell}{g_A^\ell}\right|
=\tan^2\theta_{\ell R}\lesssim1.01\times10^{-4},
\qquad \ell=\mu,\tau.
\end{gathered}
\]
The signed vector couplings have opposite overall signs in the muon and tau sectors because their \(U(1)_X\) charges are opposite; their magnitudes obey the same bound, \(|g_V^\ell|\lesssim10^{-4}|g_X|\), for both \(\ell=\mu,\tau\). The signed axial couplings likewise have opposite overall signs, while their absolute magnitudes remain close to \(|g_X|\) in this small-mixing domain. Thus, the residual vector component is suppressed by at least four orders of magnitude relative to the axial component under the adopted condition. This is a consequence of the small-mixing assumption, not a separate experimental bound; by itself it does not establish that every observable is insensitive to the vector component. We clarify the terminology of Eq.~(27) in our response to Minor Point~4.

\medskip
\noindent {\bf Minor Point 4)}.
\begin{quote}\itshape
Conservative benchmark. Equation (27) is presented as a ``conservative benchmark.'' What does this mean? A benchmark point would be a choice of these mixing angles. A constraint is an upper bound. Eq. (27) is written as an upper bound, but is described as a benchmark. This is not simply semantic: readers would want to understand if \(10^{-2}\) is an experimental upper limit or simply a convenient choice.
\end{quote}

\noindent {\bf Reply.}
Equation~(27) specifies an upper envelope for the small-mixing domain adopted in the phenomenological analysis; it is neither a single benchmark point nor a universal experimental upper limit on all four angles. We impose the same numerical upper bound separately on \(\theta_{\ell L}\) and \(\theta_{\ell R}\), for each \(\ell=\mu,\tau\); this does not set the left- and right-handed angles equal, nor does it require the muon and tau angles to coincide. These independent bounds define a common near-axial working domain. At its upper edge, the leading mixing corrections to the weak and gauge currents scale as \(\theta^2\lesssim10^{-4}\), while the right-handed rotation gives \(|g_V^\ell/g_A^\ell|=\tan^2\theta_{\ell R}\lesssim1.01\times10^{-4}\). Thus the chosen value controls both the size of the light--heavy admixture and the residual vector component; it does not by itself prove that every observable is insensitive to finite mixing, which is assessed process by process under Major Point~4.3. For comparison, the model discussion quotes LEP leptonic-width constraints on the left-handed muon mixing at the few-percent level~\cite{ParticleDataGroup:2024cfk,ALEPH:2005ab}. The adopted \(10^{-2}\) value is chosen within this small-mixing regime, but the LEP comparison concerns \(\theta_{\mu L}\) and is not a direct bound on all four angles. Other vector-like-lepton fits depend on the electroweak representation and the observables included in the fit~\cite{Crivellin:2020ebi,Poh:2017tfo}. Accordingly, \(10^{-2}\) is a deliberately specified analysis assumption, not a global-fit result or a universal experimental limit; its relation to the conditional VLL mass scales is given in Minor Point~3.

\medskip
\noindent {\bf Minor Point 5)}.
\begin{quote}\itshape
What is the purpose of including muon collider bounds? The second half of the paper focuses on existing experimental bounds on the effective theory. It is a strange change in tone to include hypothetical muon collider bounds for future reach. The authors do not justify why a future muon collider is the way to test this model in the future and the section feels like it is tacked on. Most of the section emphasizes that the muon collider search would be a ``standard'' follow up of existing work (Ref 83) and so this section does not seem particularly novel. The authors simply focus on the forward-backward asymmetry. Why not also include polarization asymmetries? Or would there be asymmetries in the \(ee\) versus \(\tau\tau\) final states? Given the first half of the paper, though, it seems that the real ``smoking gun'' for this UV completion is the physics of the \(S\) particle which should be tested in any Higgs-factory phase of a future lepton collider. I feel that this section either adds too little and should be downplayed as an example of future tests, or otherwise should have its own full fledged exploration (which may be a separate paper).
\end{quote}


\noindent {\bf Reply.}
A high-energy muon collider is a natural complementary probe because \(A'\) couples directly to the initial-state muons and can therefore be tested in the heavier-mass regime, beyond the low-energy probes emphasized in the first part of the paper. Ref.~\cite{Huang:2021nkl} provides a useful vector \(L_\mu-L_\tau\) collider benchmark at \(\sqrt{s}=3\,\mathrm{TeV}\). The axial model, however, cannot be assessed by simply importing the vector result: the left- and right-handed couplings entering the helicity amplitudes are different.

We use the \(\tau^+\tau^-\) final state because it receives only \(s\)-channel \(\gamma,Z,A'\) exchange; the same-flavor \(\mu^+\mu^-\) final state also contains \(t\)-channel diagrams, which complicate the angular interpretation. For \(\mu^+\mu^-\to\tau^+\tau^-\), neglecting the external-lepton masses, the unpolarized angular distribution is
\[
\begin{aligned}
\frac{d\sigma_{\tau\tau}}{d\cos\theta}
=\frac{s}{128\pi}\Big[&
\left(|\mathcal A_{LL}|^2+|\mathcal A_{RR}|^2\right)(1+\cos\theta)^2\\
+{}&\left(|\mathcal A_{LR}|^2+|\mathcal A_{RL}|^2\right)(1-\cos\theta)^2\Big],
\end{aligned}
\]
where the chiral amplitudes are
\[
\mathcal A_{ij}
=
\frac{e^2}{s}
+\frac{g_i^Zg_j^Z}{s-m_Z^2+i m_Z\Gamma_Z}
+\frac{g_i^\mu g_j^\tau}{s-m_{A'}^2+i m_{A'}\Gamma_{A'}},
\qquad i,j=L,R,
\]
and the Standard Model \(Z\)-boson couplings are
\[
g_Z=\frac{e}{s_Wc_W},\qquad
g_L^Z=g_Z\left(-\frac12+s_W^2\right),\qquad
g_R^Z=g_Zs_W^2.
\]
Here \(i\) labels the initial-state muon chirality and \(j\) the final-state tau chirality. We define \(\theta\) as the angle between the incoming \(\mu^-\) and outgoing \(\tau^-\) momenta in the center-of-mass frame, so \(\cos\theta>0\) is the forward \(\tau^-\) hemisphere.
The vector and axial charge assignments are
\[
\begin{aligned}
\text{vector:}\quad&
(g_L^\mu,g_R^\mu,g_L^\tau,g_R^\tau)
=(g_X,g_X,-g_X,-g_X),\\
\text{axial:}\quad&
(g_L^\mu,g_R^\mu,g_L^\tau,g_R^\tau)
=(g_X,-g_X,-g_X,g_X).
\end{aligned}
\]
The sign reversal between the left- and right-handed couplings changes the helicity amplitudes and the relative weights of the \((1+\cos\theta)^2\) and \((1-\cos\theta)^2\) terms. Thus a total-rate reinterpretation of the vector result can obscure the axial-current effect, whereas the angular distribution retains sensitivity to the chiral structure.

For the same central angular acceptance \(|\cos\theta|<0.9\) used in the figure, the cut forward--backward asymmetry is
\[
A_{\rm FB}
=
\frac{\displaystyle
\int_0^{0.9}d\cos\theta\,\frac{d\sigma}{d\cos\theta}
-
\int_{-0.9}^{0}d\cos\theta\,\frac{d\sigma}{d\cos\theta}}
{\displaystyle
\int_{-0.9}^{0.9}d\cos\theta\,\frac{d\sigma}{d\cos\theta}}
\]
The asymmetry isolates the forward--backward odd part of the distribution and is therefore a chirality-sensitive observable that does not require polarized initial beams. The fully inclusive definition corresponds to replacing the cut value \(0.9\) by \(1\). In the benchmark shown in the paper, \(\sqrt{s}=3\,\mathrm{TeV}\), \(m_{A'}=1\,\mathrm{TeV}\), and \(g_X=0.05\), the vector benchmark changes the total rate relative to the SM prediction in \(|\cos\theta|<0.9\) by about \(5\%\), whereas the axial benchmark changes it by about \(2\%\). In the pure-axial, heavy-VLL benchmark and for \(m_{A'}\gg m_\tau\), the tree-level widths into each charged Dirac flavor and each active chiral-neutrino flavor are, respectively, \(\Gamma(A'\to\ell^+\ell^-)=g_X^2m_{A'}/(12\pi)\) for \(\ell=\mu,\tau\), and \(\Gamma(A'\to\nu_\ell\bar\nu_\ell)=g_X^2m_{A'}/(24\pi)\) for \(\ell=\mu,\tau\). Summing these four open channels, while neglecting the loop-induced electron mode and with the heavy VLL pair channels kinematically closed, gives \(\Gamma_{A'}=g_X^2m_{A'}/(4\pi)\). Thus the nonresonant \(m_{A'}=1\,\mathrm{TeV}\), \(g_X=0.05\) point in the left panel has \(\Gamma_{A'}\simeq0.199\,\mathrm{GeV}\), whereas the quoted \(0.60\,\mathrm{GeV}\) width belongs to the resonance point \(m_{A'}=\sqrt{s}=3\,\mathrm{TeV}\) in the right-panel scan. In contrast, the axial benchmark produces \(\Delta A_{\rm FB}\simeq-2.9\times10^{-2}\), larger in magnitude than the vector shift for the same benchmark. This comparison demonstrates why an angular observable is more discriminating than a simple total-rate rescaling for the axial interaction.
Writing
\[
S_+=|\mathcal A_{LL}|^2+|\mathcal A_{RR}|^2,\qquad
S_-=|\mathcal A_{LR}|^2+|\mathcal A_{RL}|^2,
\]
the angular integration gives
\[
A_{\rm FB}(c_{\max})
=
\frac{c_{\max}}{1+c_{\max}^2/3}
\frac{S_+-S_-}{S_++S_-}.
\]
For the plotted cut \(c_{\max}=0.9\), the purely kinematic prefactor is approximately \(0.709\). Thus \(A_{\rm FB}\) directly compares the same-chirality and opposite-chirality helicity combinations.

Polarization observables would provide complementary information because they project separately onto initial- and final-state chiral amplitudes. Following the convention used in Ref.~\cite{Huang:2021nkl}, an initial-state left--right asymmetry can be defined with a polarized \(\mu^-\) beam and an unpolarized \(\mu^+\) beam:
\[
A_{\rm LR}^{\mu}
=
\frac{\sigma(\mu^-_L\mu^+\to\tau^+\tau^-)-\sigma(\mu^-_R\mu^+\to\tau^+\tau^-)}
{\sigma(\mu^-_L\mu^+\to\tau^+\tau^-)+\sigma(\mu^-_R\mu^+\to\tau^+\tau^-)}.
\]
With unpolarized initial beams, final-state tau polarimetry can instead measure
\[
P_\tau
=
\frac{\sigma_{\tau_R}-\sigma_{\tau_L}}
{\sigma_{\tau_R}+\sigma_{\tau_L}},
\]
where
\[
\sigma_{\tau_R}\equiv\sigma(\mu^+\mu^-\to\tau^-_R\tau^+_L),\qquad
\sigma_{\tau_L}\equiv\sigma(\mu^+\mu^-\to\tau^-_L\tau^+_R)
\]
in the relativistic limit. For the same symmetric angular acceptance \(|\cos\theta|<c_{\max}\) and in the massless-lepton approximation used above, the common angular integrals cancel in the ratios. With the first index of \(\mathcal A_{ij}\) denoting the initial \(\mu^-\) chirality and the second the final \(\tau^-\) chirality, these definitions become
\[
\begin{aligned}
A_{\rm LR}^{\mu}
&=
\frac{|\mathcal A_{LL}|^2+|\mathcal A_{LR}|^2
      -|\mathcal A_{RL}|^2-|\mathcal A_{RR}|^2}
     {|\mathcal A_{LL}|^2+|\mathcal A_{LR}|^2
      +|\mathcal A_{RL}|^2+|\mathcal A_{RR}|^2},\\[4pt]
P_\tau
&=
\frac{|\mathcal A_{LR}|^2+|\mathcal A_{RR}|^2
      -|\mathcal A_{LL}|^2-|\mathcal A_{RL}|^2}
     {|\mathcal A_{LL}|^2+|\mathcal A_{LR}|^2
      +|\mathcal A_{RL}|^2+|\mathcal A_{RR}|^2}.
\end{aligned}
\]
The first expression separates the two initial-muon chiralities, while the second separates the final tau helicities. They therefore supply information complementary to \(A_{\rm FB}\), which instead compares same- and opposite-chirality amplitude combinations after summing over both beams and final-state helicities. We retain \(A_{\rm FB}\) as the representative observable because the collider illustration in the manuscript uses unpolarized beams and does not model beam-polarization or tau-polarimetry efficiencies, analyzing powers, or systematic uncertainties. The expressions above establish the parton-level chiral information, but a numerical sensitivity or reach for either polarization observable would require those additional collider inputs and is not claimed.



The \(ee\) final state is not an equivalent tree-level channel in the pure-axial effective-current limit, because \(A'\) has no tree-level coupling to electrons. An electron coupling can arise through loop-induced kinetic mixing. The equal-mixing benchmark gives \(|\epsilon_{\gamma A'}|\leq2.12\times10^{-7}\) for \(g_X\leq0.1\) in the tabulated low-mass, on-shell range, corresponding to \(e|\epsilon_{\gamma A'}|\lesssim6.4\times10^{-8}\). This number is not a prediction for a multi-TeV muon collider: the mixing is momentum dependent, and the relevant virtuality for \(\mu^+\mu^-\to e^+e^-\) is \(q^2=s\). A collider-level \(ee\) prediction would require evaluating \(\epsilon_{\gamma A'}(s)\) with the corresponding matching condition. At the amplitude level, the direct \(\tau^+\tau^-\) channel contains the tree-level product \(g_i^\mu g_j^\tau\), whereas the \(ee\) channel contains an electron vertex proportional to \(e\epsilon_{\gamma A'}(s)\). Its rate and any associated asymmetry therefore depend on the high-energy, momentum-dependent mixing, not on the low-mass on-shell number quoted under Major Point~4.2. Thus, unlike the direct tree-level \(\tau^+\tau^-\) channel, the \(ee\) channel is loop induced in the pure-axial limit and cannot be compared quantitatively without a separate evaluation at \(q^2=s\).

We agree that \(S\) is an important complementary UV target for a future Higgs factory. The renormalizable portal operator \((H^\dagger H)(S^\dagger S)\) is allowed by the gauge symmetry; a nonzero renormalized portal coefficient induces Higgs--singlet mixing. For \(m_s<\sqrt{s_{ee}}-m_Z\), this mixing permits associated production \(e^+e^-\to Zs\). In the single-doublet/singlet mixing description, the tree-level rate is related to the Standard Model Higgsstrahlung rate by
\[
\sigma(e^+e^-\to Zs)
=\sin^2\theta_{hs}\,
\sigma_{\rm SM}(e^+e^-\to ZH)\big|_{m_H=m_s},
\]
so the production rate is controlled by the doublet component of \(s\) and the available phase space. Separately, if \(m_s<m_h/2\), the exotic decay \(h\to ss\) can be kinematically open. Its rate is controlled by the trilinear coupling \(g_{hss}\), defined by \(\mathcal L\supset-\tfrac12g_{hss}h s^2\), and is not fixed by \(\sin^2\theta_{hs}\) alone:
\[
\Gamma(h\to ss)=\frac{g_{hss}^2}{32\pi m_h}
\sqrt{1-\frac{4m_s^2}{m_h^2}}.
\]
Both \(g_{hss}\) and \(\theta_{hs}\) depend on the scalar potential and portal parameter. For the same small-portal matching benchmark specified under Major Point~1.3, the negligible-mixing description applies away from the narrow near-degenerate region identified there; that region is not covered by the small-mixing approximation. We therefore identify \(e^+e^-\to Zs\) and \(h\to ss\) as conditional complementary channels, without assigning a model-specific rate or reach or asserting negligible mixing throughout the scalar-mass range. The scalar decay pattern, including \(s\to A'A'\) when allowed, further determines the visible final states. The Higgs-factory scalar programme is therefore complementary to, rather than replaced by, the muon-collider angular test of the \(A'\) current.

Accordingly, the muon-collider result is presented as a concrete example of a future chirality-sensitive test, not as a detector-level discovery or exclusion projection. Its purpose is to show that the axial interaction has a distinctive angular signature that is not captured by a vector-to-axial total-rate reinterpretation, while the scalar \(S\) remains a separate target for a future Higgs factory.

\medskip
\noindent {\bf Minor Point 6)}. \begin{quote}\itshape
Terminology. The following points could be improved for clarity.
\end{quote}

\medskip
\noindent {\bf Minor Point 6.1)}. \begin{quote}\itshape
The word ``light'' on page 2 has no reference scale.
\end{quote}

\noindent {\bf Reply.}
The adjective "light" is not used as a quantitative mass classification. The Abstract states the broader MeV-to-TeV range, while the explicit comparison window in Fig.~2 is
\[
10^{-3}\,\mathrm{GeV}\leq m_{A'}\leq10^2\,\mathrm{GeV},
\qquad
10^{-6}\leq g_X\leq1.
\]
These numerical ranges, together with the observables included in the analysis, define the phenomenological domain; the word "light" does not introduce an additional universal mass threshold. This distinction matters because the paper combines probes with different kinematic thresholds and sensitivities over several orders of magnitude; no single mass cut separates all of them into a common \emph{light} and \emph{heavy} regime.

\medskip
\noindent {\bf Minor Point 6.2)}. \begin{quote}\itshape
At the bottom of page 11, when referring to Table II, the authors write ``representative bounds.'' What does representative mean? Does it mean these are the most stringent bounds? Are these ``some'' bounds? Is the list ``complete'' in any sense?
\end{quote}

\noindent {\bf Reply.}
Here, "representative" means selected, non-exhaustive, and source-dependent. The entries are selected to illustrate distinct scalar-mass regimes and physical mechanisms; they are neither a complete inventory of the literature nor necessarily the strongest constraint at every mass. Their numerical values and mass intervals retain the coupling normalization and source assumptions of the corresponding analyses, including the relevant production and free-streaming/trapping conditions. Table~II also distinguishes direct scalar--muon Yukawa limits from loop-induced fifth-force and photon-mediated results, whose translation requires the matching described under Major Point~1~\cite{Liu:2025zuz,Fiorillo:2026wso,Caputo:2021rux,Hoskins:1985tn,Kapner:2006si,Blakemore:2021zna}. The entries are not ranked as though they were measurements of one identical observable: they inherit different mass intervals, source environments, coupling-generation mechanisms, confidence conventions, and production or trapping assumptions. A common UV interpretation therefore requires the source-dependent matching rather than selecting the numerically smallest entry and treating it as a universal bound. The entries are not combined into a global exclusion of the UV parameter space.

\medskip
\noindent {\bf Minor Point 6.3)}. \begin{quote}\itshape
Between page 12--13: ``...do not provide a universal \(g_X\) bound, but they show that the intermediate scalar-mass region cannot simply be assumed unconstrained.'' What does intermediate mass mean? It is also not clear what the intended content of this statement is: are the authors presenting a bound? Why would the reader assume that a new particle in this range would be unconstrained?
\end{quote}

\noindent {\bf Reply.}
The intermediate-mass range refers specifically to the direct BaBar and Belle searches for
\[
e^+e^-\to\tau^+\tau^-s,\qquad
s\to e^+e^-\ \text{or}\ \mu^+\mu^-,
\]
which cover approximately \(0.04\,\mathrm{GeV}\lesssim m_s\lesssim6.5\text{--}7\,\mathrm{GeV}\)~\cite{BaBar:2020jma,Belle:2022gbl}. These searches provide direct experimental coverage of the intermediate-mass scalar region. For the present minimal UV benchmark with negligible portal mixing, the tree-level radial-scalar Yukawa couplings are \(y_{s\ell\ell}=m_\ell/v_s\) for \(\ell=\mu,\tau\). Accordingly, the UV translation uses the dimuon channel above threshold; the generic low-mass \(s\to e^+e^-\) search channel is not treated as a direct constraint on this benchmark. In the mass-proportional benchmark, the Belle limit \(\xi^{\rm lim}=0.483\) at \(m_s=1\,\mathrm{GeV}\)~\cite{Belle:2022gbl} gives
\[
v_s>\frac{v_h}{0.483}\simeq509\,\mathrm{GeV},
\qquad
y_{s\mu\mu}<2.1\times10^{-4},
\]
and, using \(m_{A'}=2g_Xv_s\),
\[
g_X<9.8\times10^{-4}
\left(\frac{m_{A'}}{\mathrm{GeV}}\right)
\]
at \(90\%\) confidence level for the illustrative \(m_s=1\,\mathrm{GeV}\) point with \(m_{A'}>0.5\,\mathrm{GeV}\), where \(s\to A'A'\) is closed. This is separate from the \(m_s\gtrsim10\,\mathrm{GeV}\) reference benchmark used for Fig.~2 and is not a universal bound on \(g_X\); for \(m_{A'}<0.5\,\mathrm{GeV}\), the visible branching fraction must be treated as described under Major Point~1.

The Belle~II references cover two distinct analyses and should not be described collectively as spin-one-only searches. The \(e^+e^-\to\mu^+\mu^-X\), \(X\to\tau^+\tau^-\) analysis covers \(3.6\text{--}10\,\mathrm{GeV}\) and explicitly includes a mass-proportional spin-zero leptophilic scalar~\cite{Belle-II:2023ydz}; the four-muon analysis covers \(0.212\text{--}9\,\mathrm{GeV}\) and also tests a muonphilic scalar~\cite{Belle-II:2024wtd}. The former has the same charged-lepton Yukawa scaling as the present radial scalar, while the latter is a complementary muonphilic benchmark. For the UV completion considered here, the additional \(s\to A'A'\) width modifies the visible dilepton or ditau branching fractions when the channel is open. The Belle~II scalar results therefore provide relevant published scalar sensitivity, but their coupling and \(\sigma\times\operatorname{BR}\) limits translate into the present UV parameters only after using the model-specific total width and visible branching fractions, as detailed under Major Point~1.

\medskip
\noindent {\bf Minor Point 6.4)}. \begin{quote}\itshape
Below (34): ``we impose Eq. (34) point by point'', it's not clear what is meant by ``point.''
\end{quote}

\noindent {\bf Reply.}
A ``point'' denotes one grid point \((m_{A'},g_X)\) in the parameter scan shown in Fig.~2. At each point, the symmetry-breaking scale is derived from
\[
v_s=\frac{m_{A'}}{2g_X},
\]
rather than varied independently. For a specified scalar mass \(m_s\), the scalar self-coupling is then determined by
\[
\lambda_S
=\frac{m_s^2}{2v_s^2}
=\frac{2g_X^2m_s^2}{m_{A'}^2}.
\]
The condition
\[
\lambda_S\leq\lambda_S^{\max}
\]
is consequently checked at each scan point for the adopted scalar benchmark. Equivalently, the dimensionless ratio
\[
\frac{\lambda_S}{4\pi}
=\frac{m_s^2}{8\pi v_s^2}
\]
must not exceed unity when the illustrative choice \(\lambda_S^{\max}=4\pi\) is used. This makes explicit that changing \(m_{A'}\) or \(g_X\) changes the derived \(v_s\), and hence changes the scalar self-coupling tested at that grid point.

Equation~(34) is checked point by point to test the theoretical consistency of the chosen scalar benchmark; it is not itself an experimental exclusion. The experimental constraints are a separate step.

\medskip
\noindent {\bf Minor Point 6.5)}. \begin{quote}\itshape
Page 21, first paragraph of Section E. What is the ``parameter region of interest''? This wasn't clearly defined anywhere.
\end{quote}

\noindent {\bf Reply.}
We agree that the phrase ``parameter region of interest'' did not specify the mass coverage of the searches considered in Section~III.E. The relevant four-muon searches cover approximately
\[
\begin{aligned}
\text{BaBar}:&\quad 0.212\text{--}10\,\mathrm{GeV},\\
\text{Belle II}:&\quad 0.212\text{--}9\,\mathrm{GeV},\\
\text{CMS}:&\quad 5\text{--}70\,\mathrm{GeV}.
\end{aligned}
\]
These intervals describe the mass coverage of the published analyses; they do not imply that the entire intervals are excluded in our model.

The relevant processes are
\[
e^+e^-\to\mu^+\mu^-A',
\qquad A'\to\mu^+\mu^-,
\]
and
\[
Z\to\mu^+\mu^-A',
\qquad A'\to\mu^+\mu^-.
\]
The reinterpretation accounts for the production rate and \(\mathrm{Br}(A'\to\mu^+\mu^-)\). The experimental analyses also impose prompt-resonance selection requirements. Their limits are applicable to our model only where the predicted \(A'\) decay length is compatible with those requirements; the mass coverage alone does not establish this condition throughout the parameter space.

In the original one-VEV model, the massless active neutrinos couple through
\[
\mathcal L\supset
g_X A'_\alpha\left(
\bar\nu_\mu\gamma^\alpha P_L\nu_\mu
-\bar\nu_\tau\gamma^\alpha P_L\nu_\tau
\right).
\]
A single massless chiral-neutrino channel contributes \(g_X^2m_{A'}/(24\pi)\); summing the two active flavors therefore gives
\[
\Gamma_{\nu\bar\nu}
=\frac{g_X^2m_{A'}}{12\pi}.
\]
Because this is a positive partial width, \(\Gamma_{\rm tot}\geq\Gamma_{\nu\bar\nu}\), and hence it gives an upper bound on the proper decay length:
\[
c\tau_{A'}=\frac{\hbar c}{\Gamma_{\rm tot}}
\leq\frac{\hbar c}{\Gamma_{\nu\bar\nu}}.
\]
At \(m_{A'}=0.212\,\mathrm{GeV}\) and \(g_X=7\times10^{-4}\), this yields \(\Gamma_{\nu\bar\nu}=2.76\times10^{-9}\,\mathrm{GeV}\) and \(c\tau_{A'}\leq0.0716\,\mu\mathrm{m}\). For a representative boost \(\gamma\simeq25\), the corresponding laboratory mean flight distance is at most \(1.79\,\mu\mathrm{m}\). This is the quantity relevant to the prompt-resonance selection; the estimate is a local benchmark check, not a statement about all masses and couplings.

The benchmark lies only about \(0.683\,\mathrm{MeV}\) above the dimuon threshold, since \(2m_\mu\simeq0.211317\,\mathrm{GeV}\). The axial dimuon width contains the threshold factor
\[
\Gamma(A'\to\mu^+\mu^-)
\simeq\frac{g_X^2m_{A'}}{12\pi}
\left(1-\frac{4m_\mu^2}{m_{A'}^2}\right)^{3/2},
\]
in the pure-axial limit. At this point the factor is approximately \(5.16\times10^{-4}\), so the visible dimuon channel is strongly phase-space suppressed. The neutrino-only estimate is consequently a conservative lifetime upper bound near threshold; any additional open channel can only shorten the lifetime. The actual prompt requirement remains analysis-dependent, and the single benchmark does not establish promptness throughout the full parameter space.

The vector-to-axial rate-matching approximation is controlled for \(m_{A'}\gg m_\mu\). 

Near the dimuon threshold, a rate-level comparison alone is insufficient to infer a detector-level exclusion.

\medskip
\noindent {\bf Minor Point 6.6)}. \begin{quote}\itshape
Page 21, end of section D: What's the significance of the \(v_s=2\,\mathrm{GeV}\) line? What does it mean that \(y_{s\mu\mu}\) ``becomes large'' (nonperturbative)? Why plot it? Is it a perturbativity bound?
\end{quote}

\noindent {\bf Reply.}
The line at \(v_s\simeq2\,\mathrm{GeV}\) follows from the scalar self-coupling condition for the reference benchmark \(m_s=10\,\mathrm{GeV}\) and \(\lambda_S^{\max}=4\pi\):
\[
v_s\geq\frac{m_s}{\sqrt{2\lambda_S^{\max}}}
=\frac{10\,\mathrm{GeV}}{\sqrt{8\pi}}
\simeq1.995\,\mathrm{GeV}.
\]
At this reference value, \(y_{s\mu\mu}=m_\mu/v_s\simeq0.053\), so \(y_{s\mu\mu}^2/(4\pi)\simeq2.2\times10^{-4}\ll1\). More generally, for the adopted \(\lambda_S^{\max}=4\pi\) criterion the boundary scales as
\[
v_s^{\min}(m_s)=\frac{m_s}{\sqrt{8\pi}}
\simeq0.1995\,m_s.
\]
Thus the plotted \(v_s\simeq2\,\mathrm{GeV}\) line corresponds specifically to the reference value \(m_s=10\,\mathrm{GeV}\); choosing a different scalar mass moves the boundary linearly with \(m_s\). The line is shown to identify where that reference scalar benchmark first satisfies the scalar-quartic perturbativity condition. It is therefore a scalar-quartic consistency boundary for the stated \(m_s\) benchmark, not a Yukawa nonperturbativity limit or an experimental exclusion. Its relation to the point-by-point scalar-mass condition is given in Major Point~1.2 and Major Point~2.

\medskip
\noindent {\large Summary:}

The paper constructs an anomaly-free axial \(U_A(1)_{L_{\mu-\tau}}\) UV completion in which the scalar \(S\) and vector-like leptons generate the charged-lepton masses. The present response makes explicit how the scalar mass, its lepton couplings, perturbativity, direct scalar searches, and its contribution to \((g-2)_\mu\) fit together, while identifying the benchmark under which Fig.~2 isolates the axial-current phenomenology.

In the dedicated response to Major Point~3, we discuss an established two-singlet type-I seesaw as a possible gauge-invariant completion of the neutrino sector: the existing charge-two scalar \(S\), an additional charge-one scalar \(\sigma\), and three right-handed neutrinos can generate a general Majorana mass matrix and realistic mixing. This is a response to the referee's neutrino-oscillation concern, not a change to the field content of the original paper. The original \(S+\mathrm{VLL}\) one-VEV model remains the baseline for the scalar, axial-current, kinetic-mixing, finite-vector, and stellar-cooling discussions above. The mass-texture argument is not a global UV parameter-space analysis.

For the additional constraints, the response gives a momentum-dependent kinetic-mixing estimate and compares the leading finite-vector effects in trident production and \((g-2)_\mu\). The Table~II entries are source-dependent scalar--muon benchmarks rather than a full-UV exclusion, whereas the \(A'\) supernova discussion concerns a distinct vector-current problem: neutrino-coalescence production and invisible widths follow from the unchanged \(L_\mu-L_\tau\) neutrino couplings, while semi-Compton production and absorption depend on the axial charged-muon current; the cited vector-model curve is used only as a reference scale, not as a full axial-model exclusion. All conclusions remain tied to the explicitly stated benchmarks and do not constitute a global exclusion of the full UV parameter space.

\noindent Sincerely yours,

\noindent Rundong Fang, Jinhui Guo, Ming Li, Jia Liu, Xiao-Ping Wang, Yiheng Xiong

\bibliographystyle{JHEP}
\bibliography{overleaf-source/axial-vector-model-building-v6/references,Replies_to_Referee_3_neutrino_references}
\end{document}
